\documentclass[11pt]{article}

\DeclareUnicodeCharacter{0301}{\'e}

\usepackage[margin=1in]{geometry}
\usepackage{multirow}
\usepackage{amssymb,amsmath,amsthm,amsfonts}
\usepackage{bm}
\usepackage{textgreek}
\usepackage{mathtools}
\usepackage{enumitem}
\usepackage[numbers,comma,sort&compress]{natbib}
\usepackage{authblk}
\usepackage{graphicx}
\usepackage[font=small]{caption}
\usepackage{subcaption} % [labelformat=simple]
\usepackage{tablefootnote} 
\usepackage{float}
\usepackage[ruled,vlined,linesnumbered]{algorithm2e}
\usepackage{algorithmic}
\renewcommand{\algorithmicrequire}{\textbf{Input: }}  
\renewcommand{\algorithmicensure}{\textbf{Output: }}
\usepackage{physics}


\usepackage{footnote}
\usepackage{xcolor}
\usepackage{mathrsfs}
\usepackage{bbm}
\usepackage{makecell}
\usepackage{pbox}
\usepackage[colorlinks]{hyperref}
\hypersetup{citecolor=blue}
% \urlstyle{same}
%\captionsetup[figure]{labelfont=bf}

\newtheorem{fact}{Fact}[section] 
\newtheorem{theorem}{Theorem}
\newtheorem{definition}{Definition}
\newtheorem{lemma}{Lemma}
\newtheorem{proposition}{Proposition}
\newtheorem{example}{Example}
\newtheorem{corollary}{Corollary}
\newtheorem{remark}{Remark}
\newtheorem{assumption}{Assumption}
\newcommand{\eqn}[1]{(\ref{eqn:#1})}
\newcommand{\eq}[1]{(\ref{eq:#1})}
\newcommand{\prb}[1]{(\ref{prb:#1})}
\newcommand{\thm}[1]{\hyperref[thm:#1]{Theorem~\ref*{thm:#1}}}
\newcommand{\cor}[1]{\hyperref[cor:#1]{Corollary~\ref*{cor:#1}}}
\newcommand{\defn}[1]{\hyperref[defn:#1]{Definition~\ref*{defn:#1}}}
\newcommand{\lem}[1]{\hyperref[lem:#1]{Lemma~\ref*{lem:#1}}}
\newcommand{\prop}[1]{\hyperref[prop:#1]{Proposition~\ref*{prop:#1}}}
\newcommand{\assum}[1]{\hyperref[assum:#1]{Assumption~\ref*{assum:#1}}}
\newcommand{\fig}[1]{\hyperref[fig:#1]{Figure~\ref*{fig:#1}}}
\newcommand{\tab}[1]{\hyperref[tab:#1]{Table~\ref*{tab:#1}}}
\newcommand{\algo}[1]{\hyperref[algo:#1]{Algorithm~\ref*{algo:#1}}}
\renewcommand{\sec}[1]{\hyperref[sec:#1]{Section~\ref*{sec:#1}}}
\newcommand{\append}[1]{\hyperref[append:#1]{Appendix~\ref*{append:#1}}}
\newcommand{\fac}[1]{\hyperref[fac:#1]{Fact~\ref*{fac:#1}}}
\newcommand{\lin}[1]{\hyperref[lin:#1]{Line~\ref*{lin:#1}}}
\newcommand{\itm}[1]{\ref{itm:#1}}
\newcommand{\bdelta}{\mathbf{\Delta}}
\renewcommand{\arraystretch}{1.25}
\newcommand{\algname}[1]{\textup{\texttt{#1}}}
\renewcommand{\i}{\mathrm{i}}
\renewcommand{\d}{\mathrm{d}}
\renewcommand{\L}{\mathcal{L}}
\newcommand{\ad}{\mathrm{ad}}
\newcommand{\T}{\mathcal{T}}
\newcommand{\B}{\Big}
\renewcommand{\P}{\mathcal{P}}
\newcommand{\simu}{\mathrm{sim}}
\newcommand{\Hc}{H_{\mathrm{control}}}
\newcommand{\He}{H_{\mathrm{eff}}}
\newcommand{\A}{\mathcal{A}}
\renewcommand{\O}{\mathcal{O}}
\newcommand{\tot}{\mathrm{tot}}
\newcommand{\G}{\mathcal{G}}
\newcommand{\E}{\mathcal{E}}
\renewcommand{\comm}{\mathrm{comm}}
\newcommand{\supp}{\mathrm{supp}}
\newcommand{\eps}{\varepsilon}
\newcommand{\mc}{\mathcal}
\newcommand{\mb}{\mathbb}
\def\>{\rangle}
\def\<{\langle}
\def\trans{^{\top}}

\newcommand{\xz}[1]{{\color{cyan} XZ: #1}}
\newcommand{\tyl}[1]{{\color{blue} Tongyang: #1}}
\newcommand{\syz}[1]{{\color{teal} (Shengyu: #1)}}
\newcommand{\syzn}[1]{{\color{teal} #1}}
\newcommand{\sz}[1]{{\color{blue} Shuo: #1}}

\allowdisplaybreaks

\title{Jointly Mitigating Physical and Trotter Errors in Hamiltonian Simulation with Commutator Bounds
}
\author{}
\date{}

\begin{document}

\maketitle

Summary of the results: 
\begin{enumerate}
    \item Provide an upper bound for the truncation error of the BCH formula of $\prod_v e^{\mc L_v}$ (each $\mc L_v$ is a Lindbladian) governed by the doubly right-nested commutators of $\mc L_v$.
    \item Provide a general condition involving the doubly right-nested commutators of Hamiltonians, where the Trotter error mitigation algorithm \cite{watson2025exponentially} achieves commutator and logarithm precision scaling, and show that $k$-local Hamiltonians and the second-quantized plane-wave electronic structure Hamiltonians satisfiy the condition.
    \item Propose a Lindbladian simulation algorithm with commutator and logarithm precision scaling. 
    \item Show that under mild physical noise, the Trotter error mitigation algorithm can simultaneously mitigate the physical noise incurring a constant complexity overhead. 
\end{enumerate}
All of the algorithms use at most two ancilla qubits. 

\textcolor{red}{Yicong: In addition to the number of interpolation nodes, try to prove that the overall sample complexity of ZNE procedure via 1D extrapolation is smaller than that by 2D extrapolation.}

\section{Introduction}

\xz{I tend to think that jointly mitigate algorithmic and physical error is a minor contribution of our work. The main technique contribution of our work is an upper bound of the truncation error of the BCH formula for general Lindbladians in terms of the doubly right-nested commutators. Compared to \cite{mizuta2025commutator}, we go beyond lattice Hamiltonian to general cases. Compared to \cite{fang2025high}, we obtain fine-grained commutator bounds, and use a different way to bound the remainder to incorporate dissipative dynamic.} 
If we use this perspective, I think our abstract should contain like \emph{Exponentially Reduced Circuit Depth for Trotterization under Physical Noise with Rigorous (General) Commutator Bounds}
The storyline of intro can be:
\begin{itemize}
    \item Hamiltonian simulation is a promising problem to show quantum advantages on near-term devices, and list some applications. \xz{We need to emphasize near-term applications, which makes it natural to use Trotter and consider physical noise. }
    \item List some Hamiltonian simulation algorithms, and state that although the Trotter method has sub-optimal dependence in $t$ and $\eps$, it has some advantages like commutator bounds, no ancilla, and complex controlled gates, which makes it more practical on near-term devices. \\
    \sz{tentative: Hamiltonian simulation is one of the most promising candidates to demonstrate practical quantum advantage on near-term devices, with applications ranging from quantum field theory to quantum chemistry and condensed matter physics.

The first explicit quantum simulation algorithm was proposed by Lloyd using product formulas, and then numerous advanced techniques have been extensively developed, including linear combination of unitaries (LCU), quantum signal processing (QSP), and qubitization. Despite the inferior cost dependence in time $t$ and precision $\epsilon$, the product formula enjoys commutator scaling and does not require ancillary qubits or complex controlled gates, enabling its implementation on near-term quantum devices. }
    \item \citet{watson2025exponentially} exponentially reduced the error dependence of Trotter using extrapolation. However, providing fully rigorous commutator bounds for general Hamiltonians is challenging. A key technical difficulty is that the convergence of the BCH formula, which the analysis relies on, is not guaranteed in all regimes (e.g., for general $k$-local Hamiltonians), making the derivation of commutator bounds difficult. Important progress was recently made by \citet{mizuta2025commutator}, who provided a rigorous commutator bound for the specific case of lattice Hamiltonians by using a subsystem Trotterization technique. This naturally leaves open the crucial question: \begin{align*}
        \textit{Can we establish a commutator bound for Trotter with extrapolation for general Hamiltonians? }
    \end{align*}  Then we can state our first contribution: \emph{We develop a new analysis, centered on doubly right-nested commutators, that provides a rigorous error bound for general Hamiltonians.} We apply our result to $k$-local Hamiltonian on lattice and second quantized Hamiltonian. The main technique ingredient is an upper bound of the truncation error of the BCH formula for general Lindbladians in terms of the doubly right-nested commutators.
    \item Besides the Trotter error, on near-term devices, physical noise is a more important source of error. The physical error can be reduced by qec or qem. qec is scalable but expensive, qem could incur an exponential sampling overhead in the circuit depth, but is more friendly to near-term devices. Introduce QEM methods, especially ZNE. 
    \item Most existing methods handle these two parts of error separately. Recent works show that the Trotter error and physical error can be simultaneously mitigated by a one-dimensional extrapolation. However, they only consider the first-order Trotter, model the physical noise by a global depolarizing channel approximately, and commutator bounds are missing. Then we ask the second question \begin{align*}
        \textit{Can we jointly mitigate physical and Trotter error while maintaining the commutator bound?}
    \end{align*}Our analysis of the truncation error of the BCH formula naturally applies to Lindbladian, which allows us to bound the extrapolation error of jointly mitigating dissipative physical error and algorithmic error in terms of the nested commutators of the Hamiltonians and noise Lindbladians. Then we can state our second contribution: \emph{We propose an algorithm that jointly mitigates physical and Trotter error with commutator bounds and provides the optimal choice of extrapolation parameters for some important Hamiltonians.} \xz{maybe add in discussion: Our analysis takes the locality of noise into account, which can be used to give a specific analysis of quantum devices with certain connectivity. }
\end{itemize}



\section{Notations}
We denote the set of all positive integers by $\mb{Z}_+$. For a sequence of operators $\A_1, \ldots, \A_{\Gamma}$, we denote $\prod_{\gamma=1}^\Gamma \A_{\gamma}= A_{\Gamma}\A_{\Gamma-1} \ldots \A_1$ and $\prod_{\gamma=\Gamma}^1 \A_{\gamma}= A_{1}\A_{2} \ldots \A_\Gamma$. We denote the right-nested commutator $[X_1, [X_2, [\cdots[X_{n-1}, X_n]\cdots]]]$ by $[X_1, X_2, \ldots, X_n]$. We denote $\{1, 2, \ldots, n\}$ by $[n]$ and the collection of all size-$m$ subsets of $[n]$ by $\binom{[n]}{m}$. We use $\supp(X)$ to denote the set of qubits that $X$ nontrivially acts on.  For a superoperator $\mc L$, we denote the diamond norm of $\mc L$ by $\|\mc L\|$.
Define the time-ordering operator $\mc T\{\cdot\}$ as \begin{align*}
    \mc T\big\{[O_1(t_1),O_2(t_2),\ldots, O_q(t_q) ]\big\} = [O_{\pi(1)}(t_{\pi(1)}),O_{\pi(2)}(t_{\pi(2)}),\ldots, O_{\pi(q)}(t_{\pi(q)})],
\end{align*}
where $\pi\colon [q]\to[q]$ is a permutation of $[q]$ such that \[
    t_{\pi(1)}\ge t_{\pi(2)}\ge \cdots \ge t_{\pi(q)}.
    \]
For any permutation $\sigma\colon [q]\to [q]$, define the reordering operator $\mc R_{\sigma}\{\cdot\}$ as \begin{align*}
    \mc R_{\sigma}\big\{[O_1(t_1),O_2(t_2),\ldots, O_q(t_q)]\big\} = [O_{\sigma(1)}(t_{\sigma(1)}),O_{\sigma(2)}(t_{\sigma(2)}),\ldots, O_{\sigma(q)}(t_{\sigma(q)})].
\end{align*}
Throughout this paper, we denote the Lindblad dissipator with a jump operator $L$ as $\mathcal{D}[L](\rho) := L\rho L^\dagger - \frac{1}{2}\{L^\dagger L, \rho\}$.

\section{Preliminaries}
\subsection{Quantum error mitigation}
\textit{Quantum error mitigation} (QEM)~\cite{cai2023quantum} refers to algorithmic schemes that reduce the noise-induced bias in the observable expectation value by post-processing outputs from an ensemble of circuit runs. QEM has enabled the quantum hardware in the noisy intermediate-scale quantum era to achieve practical applications \sz{to tweak: demonstrate practical quantum advantages}.

One of the most common QEM techniques is \textit{zero-noise extrapolation} (ZNE), which tries to obtain the zero-noise value using data points at different circuit fault rates. Specifically,

\subsection{Magnus expansion}
Consider a time-dependent linear differential equation \begin{align*}
    \frac{\d Y(t)}{\d t} = A(t) Y(t), \quad Y(0) = I.
\end{align*} The Magus expansion expresses $Y(t)$ as the exponential of a time-dependent operator $Y(t) = \exp(\Omega(t))$, where the exponent $\Omega(t)$ admits a series expansion \begin{align*}
    \Omega(t) = \sum_{q = 1}^{\infty} \Omega_q(t).
\end{align*} 
The $p$-th order Magnus expansion is denoted by \begin{align}
    \Omega_{(p)} = \sum_{q=1}^{p} \Omega_q. 
\end{align}
Each $\Omega_q(t)$ can be written as an integral of nested commutators of $A(t)$:
\begin{align}
\label{eq:magnus-formula}
    \Omega_q(t) = \sum_{\sigma\in S_q} c_{\sigma, q}  \int_{0}^t \d \tau_1 \int_{0}^{\tau_1}\d \tau_2 \cdots \int_{0}^{\tau_{q-1}}\d \tau_q \, \mc R_{\sigma}\big\{[A(\tau_1), A(\tau_2), \ldots, A(\tau_q)] \big\},
\end{align}
where $S_q$ denotes the set of permutations of $[q]$, \begin{align}
\label{eq:def-c-sigma}
    c_{\sigma, q}=\frac{1}{q^2}\frac{(-1)^{d_{\sigma}}}{\binom{q-1}{d_\sigma}} 
\end{align} satisfying $|c_{\sigma, q}|\le 1/q^2$, and $d_{\sigma}$ denotes the number of $i\in[q-1]$ such that $\sigma(i)>\sigma(i+1)$. The derivative of $\Omega_q(t)$ is
\begin{align}
    \dot{\Omega}_q(t) = \sum_{\sigma\in S_q} c_{\sigma, q}  \int_{0}^{t}\d \tau_2 \int_0^{\tau_2} \d\tau_3 \cdots \int_{0}^{\tau_{q-1}}\d \tau_q \, \mc R_{\sigma}\big\{[A(t), A(\tau_2), \ldots, A(\tau_q)] \big\}.
\end{align}

\subsection{Richardson extrapolation}

\begin{lemma}[{see \cite[Lemma A.9]{wang2025randomized}}]
\label{lem:richardson-extrapolation}
    Suppose that function $f\colon \mb R \to \mathbb{R}$ can be decomposed as 
    \begin{align*}
        f(x) = P_m(x)+R_m(x),
    \end{align*} 
    where $P_m(x)$ is a degree-$(m-1)$ polynomial and $R_m(x)$ is the remainder. Given any $t > 0$ and $s\in (0,1)$, we can choose 
    \begin{align*}
        K = \max\Big\{\frac{m}{\pi}, \frac{2t}{s}\Big\}, \quad r_j =\left\lceil \frac{K}{\sin^2(\pi(2j-1)/8m)}\right\rceil,
    \end{align*}
    and
    \begin{align*}
        s_j = \frac{t}{r_j}, \quad b_j=\prod_{\ell \neq j} \frac{1}{1-r_\ell/r_j}
    \end{align*}
    so that $1\le \|\mathbf{b}\|_1 \le C\log m$ for some absolute constant $C$, $r_j=O(\max\{m^3, m^2t/s\}/j^2)$, and $\max_{j}\{ s_j \}\le s$. Then define the $m$-term Richardson extrapolation \begin{align}
        \label{eq_richardson extrapolation}
        F^{(m)}(s)=\sum_{j=1}^m b_j f(s_j),
    \end{align}
    which satisfies \begin{align}
    \label{eq:error-F}
        |F^{(m)}(s)-f(0)|\le \|\mathbf{b}\|_1 \max_j|R_{m}(s_j)|.
    \end{align}
\end{lemma}
\subsection{Technical lemmas}
Let $\mc A(\tau)$ and $\mc B(\tau)$ be two continuous operator-valued functions. We rely on the following lemma to decompose the time-ordered exponential of the sum $\mc A(\tau) + \mc B(\tau)$. This result is essentially the interaction picture representation of the evolution.
\begin{lemma}[{\cite[Page 21]{dollard1984product}}]
\label{lem:interaction-pic}
    Let $\mc H(\tau) = \mc A(\tau) + \mc B(\tau)$ be an operator-valued function defined for $\tau \in \mathbb{R}$ with continuous summands $\mc A(\tau)$ and $\mc B(\tau)$. Then
    \begin{align*}
    \exp_{\mc T}\left(\int_{0}^{t}  \mc H(\tau)\, \d \tau\right) =  \begin{multlined}[t]
        \exp_{\mc T}\left(\int_{0}^{t}  \mc A(\tau)\, \d \tau\right) \\
     \cdot \exp_{\mc T}\left(\int_{0}^{t}\exp_{\mc T}^{-1}\left(\int_{0}^{\tau_1}  \mc A(\tau_2)\, \d \tau_2\right)\mc  B(\tau_1) \exp_{\mc T}\left(\int_{0}^{\tau_1}\mc  A(\tau_2)\, \d \tau_2\right)\, \d \tau_1\right).
    \end{multlined}
    \end{align*}
\end{lemma} 
We use the following lemma to bound the difference between powers of a superoperator and a quantum channel. 
\begin{lemma}
\label{lem:power-bound}
    Let $\mc N_1, \mc N_2$ be two superoperators, where $\mc N_1$ is a quantum channel. For any positive integer $k$, the difference between $\mc N_1^k$ and $\mc N_2^{k}$ can be bounded as
    \begin{align*}
        \|\mc N_1^k-\mc N_2^k\| \le k \max\{1, \|\mc N_2\|\}^{k-1} \|\mc N_1-\mc N_2\|.
    \end{align*}
\end{lemma}
\begin{proof}
    We use the telescoping sum identity for the difference between powers of two operators:
    \begin{align*}
        \mc N_1^k - \mc N_2^k = \sum_{j=0}^{k-1} \mc N_1^{k-1-j} (\mc N_1 - \mc N_2) \mc N_2^j.
    \end{align*}
    Taking the diamond norm on both sides, we obtain
    \begin{align*}
        \|\mc N_1^k - \mc N_2^k\| &\le \sum_{j=0}^{k-1} \|\mc N_1\|^{k-1-j} \|\mc N_1 - \mc N_2\| \|\mc N_2\|^j \\
        &= \|\mc N_1 - \mc N_2\| \sum_{j=0}^{k-1} \|\mc N_2\|^j \\
        &\le k\max\{1, \|\mc N_2\|\}^{k-1}\|\mc N_1-\mc N_2\|,
    \end{align*}
    where the second line follows from the fact that $\mc N_1$ is a quantum channel, so $\|\mc N_1\|=1$. 
\end{proof}
\section{Truncation error of the BCH formula}
\label{sec:bch-truncate}
Let $\mc K_1, \ldots, \mc K_M$ be general Lindbladians. The BCH formula can then be written as
\begin{align}
\label{eq:def-bch}
 	e^{\mc K_M}\cdots e^{\mc K_2}e^{\mc K_1} = \exp(\sum_{q = 1}^{\infty} \Phi_q),
\end{align}
where each $\Phi_q$ is a sum of the weighted right-nested commutators of $\mc K_j$ (see \citet{arnal2021note}):
\begin{align}
\label{eq:bch}
 	\Phi_q = \sum_{\substack{p_v \ge 0,\\p_1+\cdots+p_M = q}}\frac{1}{p_1!\cdots p_M!} \sum_{\sigma\in S_q} c_{\sigma, q}\mc R_{\sigma}\{[\underbrace{\mc K_M, \ldots, \mc K_M}_{p_M}, \ldots, \underbrace{\mc K_1, \ldots, \mc K_1}_{p_1}]\},
\end{align}
and $c_{\sigma,q}$ is the same coefficient as in Eq.~\eq{def-c-sigma}. Define the sum of the $q$-fold right-nested commutator norms as
\begin{align}
\label{eq:def-comm}
 	\alpha_{\rm comm}^{(q)} =\sum_{v_1, \ldots, v_q = 1}^{M} \left\|[\mc K_{v_1}, \ldots, \mc K_{v_q}]\right\|.
\end{align}
\citet{aftab2024multi} showed that the norm of $\Phi_q$ can be bounded as
\begin{align*}
 	\|\Phi_q\|\le \frac{1}{q^2} 	\alpha_{\rm comm}^{(q)},
\end{align*}
hence, the series expansion in Eq.~\eq{bch} converges if there exist two constants $J,C\ge 0$ such that
\begin{align}
\label{eq:converg-bch}
 	\sup_{q\ge J}\alpha_{\rm comm}^{(q)} \le C.
\end{align} 
A trivial bound for $\alpha_{\rm comm}^{(q)}$ is
\begin{align}
\label{eq:trivial-alpha}
 	\alpha_{\rm comm}^{(q)} \le \left(2\sum_{v = 1}^M \|\mc K_v\|\right)^q,
\end{align}
by $\|[A,B]\|\le 2\|A\|\|B\|$ and the triangle inequality. The quantity $\alpha_{\rm comm}^{(q+1)}$ also appears in the complexity of $q$-th order Trotterization \cite{childs2021theory}, where each $\mc K_v$ is an Hermitian operator in the decomposition of the simulated Hamiltonian. For some physical Hamiltonians, such as $k$-local Hamiltonians on a lattice of size $N$, where each $\mc K_v$ can be decomposed as a sum of terms acting nontrivially on at most $k$ sites, $\alpha_{\rm comm}^{(q)}$ admits a sharper bound than Eq.~\eq{trivial-alpha}. This leads to the commutator scaling of the $q$-th order Trotterization. Specifically, suppose that the maximum interaction strength of $\sum_{v = 1}^M \mc K_v$ on each site is $g$. Eq.~\eq{trivial-alpha} yields $(2Ng)^q$, while the bound considering the commutator structure of $\mc K_v$ gives $Nq!(2kg)^q = O(N(qkg)^q)$. The latter bound gives better complexity scaling for local Hamiltonians where $k\ll N$ and $q$ is a constant. 

However, in the error analysis of some advanced Trotter methods with higher precision \cite{aftab2024multi, watson2025exponentially}, bounds on $\alpha_{\rm comm}^{(q)}$ for a single $q$ do not suffice. Instead, we require the BCH formula to be convergent. The convergence condition in Eq.~\eq{converg-bch} requires $\alpha_{\rm comm}^{(q)}$ to be bounded for arbitrarily large $q$, and is thus too strict to utilize the commutator structure of $\mc K_v$, since the only valid choice to satisfy the condition is $g = 0$. \xz{add connection to trotter, and comparision to previous analysis} The situation is the same if each $\mc K_v$ is a $k$-local Lindbladian, as the parameter $\alpha_{\rm comm}^{(q)}$ admits similar bounds (see \cor{commu-bound-sum}). 

To bypass the divergence issue of the BCH formula in the analysis of the multiproduct formula, \citet{mizuta2025commutator} showed that for $k$-local Hamiltonians on a lattice, the truncated BCH formula could approximate the exponential product even outside its convergence radius. Consequently, we analyze the truncated BCH formula instead of the full series, thereby avoiding the convergence requirement Their proof relies on a subsystem Trotterization technique specific to lattices, which is difficult to extend to general Hamiltonians. To analyze the BCH truncation error for general cases, we construct an equivalent continuous-time evolution via the Magnus expansion

We first define a continuous-time differential equation constructed to match the discrete product in Eq.~\eq{def-bch}. Let\begin{align}\label{eq:df-eq-general}\frac{\d \mc Y(\tau)}{\d \tau} = \mc L(\tau) \mc Y(\tau), \quad \mc Y(0)= \mc I, \quad \mc L(\tau) = \mc K_j \quad \text{if } \tau \in (j-1, j] \text{ for } j=1,\ldots,M.\end{align}The solution to this ODE at time $\tau=M$ is exactly the discrete product in Eq.~\eq{def-bch}:\[\mc Y(M)  = \exp(\mc K_M) \cdots \exp(\mc K_1) .
\]
On the other hand, the solution $\mc Y(M)$ can also be expressed using the Magnus expansion for the generator $\mc L(\tau)$ defined in Eq.~\eq{df-eq-general}. Let $\mc Y(M) = \exp(\Omega(M))$, where $\Omega(M) = \sum_{q=1}^{\infty} \Omega_q(M)$ is the Magnus expansion.

We now show that each term in the BCH formula ($\Phi_q$) and the Magnus expansion ($\Omega_q(M)$) is identical. The proof is based on the formal expression of these terms in Eq.~\eq{def-bch} and Eq.~\eq{magnus-formula} and does not require the convergence of the BCH formula and the Magnus expansion.  
\begin{lemma}[Equivalence of BCH and Magnus Generators]
\label{lem:magnus-bch-general}
The $q$-th order term $\Phi_q$ of the BCH series in Eq.~\eq{def-bch} is identical to the $q$-th order Magnus expansion term $\Omega_q(M)$ in Eq.~\eq{magnus-formula} evaluated at $\tau=M$ for the continuous system defined in Eq.~\eq{df-eq-general}.
\end{lemma}
\begin{proof}
    By definition, the $q$-th order Magnus expansion term $\Omega_q(M)$ for the continuous system in Eq.~\eq{df-eq-general} is given by the integral formula (Eq.~\eq{magnus-formula} with $t=M$)
    \begin{align}
         \Omega_q(M) = \sum_{\sigma\in S_q} c_{\sigma, q}  \int_{0}^{M} \d \tau_1 \int_{0}^{\tau_1}\d \tau_2 \cdots \int_{0}^{\tau_{q-1}}\d \tau_q ~\mc R_{\sigma}\big\{[\mc L(\tau_1), \ldots, \mc L(\tau_q)] \big\}. \nonumber
    \end{align}
    We first symmetrize the time-ordered integral by extending the integration domain to the hypercube $[0, M]^q$ and discretize the integral, which yields
    \begin{align*}
        \Omega_q(M) =& \frac{1}{ q!}\sum_{\sigma\in S_q} c_{\sigma, q}  \int_{0}^{M} \d \tau_1 \cdots \int_{0}^{M}\d \tau_q ~ \mc R_{\sigma} \circ \mc T\big\{[\mc L(\tau_1), \ldots, \mc L(\tau_q)]\big \} \\
         =&\frac{1}{ q!}\sum_{\sigma\in S_q} c_{\sigma, q} \sum_{ 1\le v_1,\ldots, v_q \le M} \int_{v_1-1}^{v_1} \d \tau_1 \cdots \int_{v_q-1}^{v_q}\d \tau_q ~ \mc R_{\sigma} \circ \mc T\big\{[\mc L(\tau_1), \ldots, \mc L(\tau_q)] \big\}.
    \end{align*}
    Since $\mc L(\tau)$ is constant on each interval $(v_j-1, v_j]$ the integral simplifies to the sum
    \begin{align}
        \Omega_q(M) = \frac{1}{ q!}\sum_{\sigma\in S_q} c_{\sigma, q} \sum_{ 1\le v_1,\ldots, v_q \le M} \mc R_{\sigma} \circ \mc T\big\{[\mc K_{v_1}, \ldots, \mc K_{v_q}] \big\}. \label{eq:magnus-general-midstep}
    \end{align}
    Now, we regroup the inner sum over the $M^q$ sequences $(v_1, \ldots, v_q)$ according to the number of times each generator $\mc K_j$ appears. Let $p_j \ge 0$ be the count of index $j$ in a sequence, such that $p_1 + \cdots + p_M = q$. The number of sequences $(v_1, \ldots, v_q)$ corresponding to a specific set of counts $(p_1, \ldots, p_M)$ is given by the multinomial coefficient $\binom{q}{p_1, \ldots, p_M}$.
    
    Regrouping the sum in Eq.~\eq{magnus-general-midstep} by these counts yields
    \begin{align*}
         \Omega_q(M) = &\frac{1}{ q!}\sum_{\sigma\in S_q} c_{\sigma, q} \sum_{\substack{p_j \ge 0,\\ p_1+\cdots+p_M = q}} \binom{q}{p_1, \ldots, p_M} R_{\sigma}\{[\underbrace{\mc K_M, \ldots, \mc K_M}_{p_M}, \ldots, \underbrace{\mc K_1, \ldots, \mc K_1}_{p_1}]\} \\
         =&\sum_{\substack{p_j \ge 0,\\ p_1+\cdots+p_M = q}} \frac{1}{p_1!\cdots p_M!} \sum_{\sigma\in S_q} c_{\sigma, q} \mc R_{\sigma}\{[\underbrace{\mc K_M, \ldots, \mc K_M}_{p_M}, \ldots, \underbrace{\mc K_1, \ldots, \mc K_1}_{p_1}]\},
    \end{align*}
    which matches $\Phi_q$ in Eq.~\eq{def-bch}.

\end{proof}
Establishing this equivalence allows us to bound the BCH truncation error by studying the corresponding Magnus expansion. Before the analysis, we first introduce the necessary concepts using the general notation from our preliminaries.
We define the  $q_0$-th order truncated generator  as
\begin{align*}
 \Omega_{(q_0)}(\tau) := \sum_{q'=1}^{q_0} \Omega_{q'}(\tau).
\end{align*}
The corresponding $q_0$-th order approximate evolution is
\begin{align*}
 \mc Y_{(q_0)}(\tau):=\exp\left( \Omega_{(q_0)}(\tau) \right).
\end{align*}
A key result from the theory of Magnus expansions (see \cite[Section 3.2]{fang2025high}) states that this approximate evolution $\mc Y_{(q_0)}(\tau)$ is itself the exact solution of a differential equation governed by a modified generator $\mc{L}_{(q_0)}(\tau)$:
\begin{align*}
 \frac{\d \mc Y_{(q_0)}(\tau)}{\d \tau} = \mc{L}_{(q_0)}(\tau) \mc Y_{(q_0)}(\tau).
\end{align*}
This modified generator $\mc{L}_{(q_0)}(\tau)$ is given by the series expansion
\begin{align}
\label{eq:expansion-L-q}
 \mc{L}_{(q_0)}(\tau) = \sum_{\ell = 0}^{\infty}\frac{1}{(\ell+1)!}\ad_{\Omega_{(q_0)}(\tau)}^\ell(\dot{\Omega}_{(q_0)}(\tau)).
\end{align}
Our first step is to reduce the total truncation error to the integral of the generator error.
\begin{lemma}
\label{lem:error-reduction-general}
The error of the $q_0$-th order approximation evolution is bounded by
\begin{align*}
\|\mc \exp(\Omega_{(q_0)}(M)) -\mc Y(M)\| \le \bigg( \int_{0}^{M}\| \mc{L}_{(q_0)}(\tau)-\mc{L}(\tau)\|\,\d \tau\bigg)\exp\bigg(\int_{0}^{M}\| \mc{L}_{(q_0)}(\tau)-\mc{L}(\tau)\|\,\d \tau\bigg).
\end{align*}
\end{lemma}
\begin{proof}
Let $\mc Y(\tau_2, \tau_1) = \exp_{\mc T}(-\int_{\tau_1}^{\tau_2}\mc L(\tau')\,\d \tau' )$. As $\mc L(\tau)$ is a Lindbladian, $\mc Y(\tau_2, \tau_1)$ is a quantum channel and $\|\mc Y(\tau_2, \tau_1)\| = 1$. The exponential of $\mc L_{(q_0)}(\tau)$ can be bounded as
\begin{align*}
 &\bigg\|\exp_{\mc T}\bigg(\int_{0}^{\tau} \mc L_{(q_0)}(\tau')\,\d \tau'\bigg)\bigg\| \\
 \le ~& \|\mc Y(\tau,0)\|+\bigg\|\sum_{j=1}^{\infty}\int_{0}^{\tau}\d \tau_1\int_0^{\tau_1}\d \tau_2 \cdots \int_{0}^{\tau_{j-1}}\d \tau_j~\prod_{j'=j}^{1}\Big(\mc Y(\tau_{j'-1},\tau_{j'})(\mc L_{(q_0)}(\tau_{j'})-\mc L(\tau_{j'}))\Big)\mc Y(\tau_j,0)\bigg\|\\
 \le ~&1+\sum_{j=1}^{\infty}\int_{0}^{\tau}\d \tau_1\int_0^{\tau_1}\d \tau_2 \cdots \int_{0}^{\tau_{j-1}}\d \tau_j~\prod_{j'=j}^{1}\| \mc L_{(q_0)}(\tau_{j'})-\mc L(\tau_{j'})\|\\
 =~& \exp\bigg(\int_{0}^{\tau}\| \mc L_{(q_0)}(\tau')-\mc L(\tau')\|\,\d \tau'\bigg),
\end{align*}
where the second line follows from applying \lem{interaction-pic} with $\mc A(\tau) = \mc L(\tau)$ and $\mc H(\tau)= \mc L_{(q_0)}(\tau)$ and expanding the result into a Dyson series. The fourth line follows from the fact that $\mc L(\tau)$ generates a quantum channel, implying $\|\mc Y(t_2, t_1)\| = 1$ for any $t_2 \ge t_1$. 
Using Duhamel's principle, the difference between $\mc Y_{(q_0)}(M)$ and $\mc Y(M)$ as
\begin{align*}
 \|\exp(\Omega_{(q_0)}(M)) -\mc Y(M)\|=&~ \|\mc Y_{(q_0)}(M) -\mc Y(M)\|\\
 =&~\bigg\|\exp_{\mc T}\bigg(\int_{0}^{M} \mc L_{(q_0)}(\tau)\,\d \tau\bigg)-\exp_{\mc T}\bigg(\int_{0}^{M} \mc L(\tau)\,\d \tau\bigg)\bigg\| \\
 =&~\bigg\|\int_{0}^{M} \mc Y(M, \tau)(\mc L_{(q_0)}(\tau)-\mc L(\tau)) \exp_{\mc T}\bigg(\int_{0}^\tau \mc L_{(q_0)}(s)\,\d s\bigg)\,\d \tau\bigg\| \\
 \le&~\bigg( \int_{0}^{M}\| \mc L_{(q_0)}(\tau)-\mc L(\tau)\|\,\d \tau\bigg)\exp\bigg(\int_{0}^{M}\| \mc L_{(q_0)}(\tau)-\mc L(\tau)\|\,\d \tau\bigg).
\end{align*}
where the third line follows from $\|\mc Y(M,\tau)\| = 1$.
\end{proof}

\lem{error-reduction-general} reduces our main task to analyzing and bounding $\int_{0}^{M}\| \mc{L}_{(q_0)}(\tau)-\mc{L}(\tau)\|$. To understand the structure of this error term, we rely on the following properties of the modified generator, which are based on the grade of the commutator terms \cite{fang2025high}. For an expression $L$ involving nested time integrals of commutators of $\mc L(\tau)$ (or norms of such terms), its \emph{grade} $\mathrm{gd}(L)$ is defined as the total number of $\mc L$ operators appearing in its nested commutator structure.

\begin{lemma}[{\cite[Proposition~3]{fang2025high}}]
\label{lem:grade-cond}
 For any $1\le k\le q_0$, the sum of all terms in $\mc L_{(q_0)}(\tau) -\mc L(\tau)$ with grade $k$ vanishes. 
\end{lemma} 
\begin{lemma}[{\cite[Lemma~4]{fang2025high}}]
\label{lem:grade-cond-q+1}
 The grade $q_0+1$ terms in $\mc L_{(q_0)}(\tau) -\mc L(\tau)$ match those in $-\dot\Omega_{q_0+1}(\tau)$.
\end{lemma}

These lemmas show that the error integral in \lem{error-reduction-general} contains only terms with grade $q_0+1$ and higher. As highlighted in the introduction, the standard right-nested commutator bounds scale poorly (e.g., as $q!$) and lead to divergence for $k$-local systems. 

This motivates our core technical insight: we show that all higher-grade terms ($q > q_0$) in $\int_{0}^{M}\| \mc L_{(q_0)}(\tau)-\mc L(\tau)\|\,\d \tau$ are composed of \emph{doubly right-nested commutators} of $\mc K_v$. For $k$-local Hamiltonians and second-quantized plane-wave electronic structure Hamiltonians, these doubly right-nested commutators admit a better scaling than right-nested ones, allowing us to derive a meaningful bound.

Define the sum of doubly right-nested commutator norms as
\begin{align}
\label{eq:def-beta-bar-general}
    \bar{\beta}_{\comm}^{(q_1, \ldots, q_{d})}
    =~ 
    \sum_{v_{1,1}, \ldots, v_{d, q_{d}} = 1}^M \bigg\| \Big[ 
         [\mc K_{v_{1,1}}, \ldots, \mc K_{v_{1,q_1}}], \ldots, 
        [\mc K_{v_{d,1}}, \ldots, \mc K_{v_{d,q_{d}}}]
     \Big] \bigg\|.
\end{align} We now state the lemma that bounds the integrated expansion terms in $\mc{L}_{(q_0)}(\tau)$ as doubly right-nested commutators of the generators. \xz{maybe explain the what is the better scaling of the doubly right-nested commutator here for some spedific hamiltonian.}

\begin{lemma}
\label{lem:doubly-nested-bound-general}
For any sequence of positive integers $q_1, \ldots, q_{d+1}$, the integrated norm of the corresponding term in the modified generator $\mc L_{(q_0)}$ is bounded as
 \begin{align*}
  \int_{0}^{M} \bigg\|\prod_{r=d}^{1} \ad_{\Omega_{q_{r}}(\tau)}\cdot\dot\Omega_{q_{d+1}}(\tau)\bigg\| \,\d\tau \le \bigg(\prod_{r=1}^{d+1} \frac{1}{q_r^2}\bigg) q_{d+1}\bar{\beta}_{\comm}^{(q_1, \ldots, q_{d+1})}.
 \end{align*}
\end{lemma}
\begin{proof}
The proof follows a similar, though more involved procedure to the one used in the proof of \lem{magnus-bch-general}. We substitute the integral definitions for each $\Omega_q(\tau)$ and $\dot\Omega_{q+1}(\tau)$ term, apply the triangle inequality, and perform a detailed combinatorial analysis by discretizing the resulting time-ordered integrals and regrouping the resulting sums. The full derivation is deferred to \append{sec:proof-doubly-nested}.
\end{proof}



\begin{theorem}[BCH Truncation Bound]
 \label{thm:general-bch-truncation}
 For any positive integer $q_0$, define 
 \begin{align}
\label{eq:def-beta-comm-sum}
    \bar{\beta}_{\comm,q_0} := \bar{\beta}_{\comm}^{(q_0+1)} + \sum_{r = 0}^{\infty}\frac{1}{(r+1)!} \sum_{\substack{1\le p_1, \ldots, p_{r+1}\le q_0\\p_1+\cdots+p_{r+1} \ge q_0+2}}\bar{\beta}_{\comm}^{(p_1, \ldots, p_{r+1})}.
\end{align}
 If $\bar{\beta}_{\comm,q_0} \le 1$, the truncation error of the $q_0$-th order BCH formula is bounded as
 \begin{align*}
 \left\|\exp\Big(\sum_{q=1}^{q_0} \Phi_q\Big) - \prod_{v = 1}^Me^{\mc K_v} \right\| \le e\bar{\beta}_{\comm,q_0}.
 \end{align*}
\end{theorem}
\begin{proof}
    By \lem{magnus-bch-general}, the problem is equivalent to bounding the Magnus expansion truncation error $\|\exp(\Omega_{(q_0)}(M)) - \mc Y(M)\|$ for the continuous system in Eq.~\eq{df-eq-general}, where $\Omega_{(q_0)}(M)$ is the $q_0$-th order Magnus generator. 

    By \lem{error-reduction-general}, it suffices to bound $\int_{0}^{M}\| \mc L_{(q_0)}(\tau)-\mc L(\tau)\|\,\d \tau$. By \lem{grade-cond}, we only need to consider terms with a grade greater than $q_0$. By \lem{grade-cond-q+1}, the grade-$(q_0+1)$ term of $\mc L_{(q_0)}(\tau)-\mc L(\tau)$ is $-\dot\Omega_{q_0+1}(\tau)$. Applying \lem{doubly-nested-bound-general} with $d = 0$ and $q_1 = q_0+1$ yields
    \begin{align*}
     \int_{0}^{M}\|\dot\Omega_{q_0+1}(\tau)\|\, \d \tau\le\bar{\beta}_{\comm}^{(q_0+1)}.
    \end{align*}
    For any $q>q_0+1$, as $\mc L(\tau)$ is grade-one, the grade-$q$ terms of $\mc L_{(q_0)}(\tau)-\mc L(\tau)$ match the grade-$q$ terms of $\mc L_{(q_0)}(\tau)$, which are
    \begin{align*}
        \sum_{k=0}^{\infty}\frac{1}{(k+1)!}\sum_{\substack{1\le p_1, \ldots, p_{k+1}\le q_0\\p_1+\cdots+p_{k+1}=q}}\ad_{\Omega_{(q_0)}(\tau)}^k(\dot{\Omega}_{(q_0)}(\tau)),
    \end{align*} 
    according to Eq.~\eq{expansion-L-q}. Applying \lem{doubly-nested-bound-general} to the grade-$q$ terms of $\int_{0}^M \|\mc L_{(q_0)}(\tau)-\mc L(\tau)\|$ yields 
    \begin{align*}
     \sum_{k=0}^{\infty}\frac{1}{(k+1)!}\sum_{\substack{1\le p_1, \ldots, p_{k+1}\le q_0\\p_1+\cdots+p_{k+1}=q}}\bar{\beta}_{\comm}^{(p_1, \ldots, p_{k+1})}.
    \end{align*}
    Therefore, the difference integral can be bounded by
    \begin{align*}
     \int_{0}^{M}\|\mc L_{(q_0)}(\tau)-\mc L(\tau)\|\,\d \tau\le \bar{\beta}_{\comm}^{(q_0+1)}+ \sum_{q = q_0+2}^{\infty} \sum_{k=0}^{\infty}\frac{1}{(k+1)!}\sum_{\substack{1\le p_1, \ldots, p_{k+1}\le q_0\\p_1+\cdots+p_{k+1}=q}}\bar{\beta}_{\comm}^{(p_1, \ldots, p_{k+1})} = \bar\beta_{\mathrm{comm}, q_0}.
    \end{align*}
    Therefore, the error between the two exponentials can be bounded as
    \begin{align*}
    \|\exp(\Omega_{(q_0)}(M))-\mc Y(M)\|&\le \bigg( \int_{0}^{M}\| \mc L_{(q_0)}(\tau)-\mc L(\tau)\|\,\d \tau\bigg)\exp\bigg(\int_{0}^{M}\| \mc L_{(q_0)}(\tau)-\mc L(\tau)\|\,\d \tau\bigg)\\
    &\le e^{\bar\beta_{\mathrm{comm}, q_0}}\bar\beta_{\mathrm{comm}, q_0} \le e\bar\beta_{\mathrm{comm}, q_0},
    \end{align*}
    which concludes the proof.
\end{proof}
\xz{TBD: explain $\bar\beta_{\mathrm{comm}, q_0}$, it consider only right-nested commutator with fixed grade $q_0+1$, and the other terms are right-nested commutator of at most $q_0$-fold right-nested commutator, where the total grade is larger than $q_0 + 1$.}

%%%%%%%%%%%%%%%%%%%%%%%%%%%%%%%%%%%%%%%%%%%%%%%%%%%%%%%%%%%%%%%%%%%%%%%%%%%%%%%%%%%%%%%%%%%%%%%%%%%%%%%%%%%%%
\subsection{Application to $k$-local systems}
\label{sec:k-local-example}

We now demonstrate the power of \thm{general-bch-truncation} by applying it to the $k$-local systems that motivated our analysis. Suppose that each generator $\mc K_v$ is a \emph{$k$-local superoperator} on a lattice $\Lambda = [N]$. Formally, this means
\begin{align*}
    \mc K_v = \sum_{\gamma = 1}^{\Gamma} \mc K_{\gamma}^{(v)}, \text{ with } |\supp(\mc K_{\gamma}^{(v)} )|\le k, \quad \forall v\in [M], \gamma\in [\Gamma].
\end{align*} 
Note that while our motivation comes from Lindbladians, the following derivation relies solely on the locality structure and holds for general superoperators.
For any site $j\in \Lambda$, we further suppose that 
\begin{align*}
    \sum_{v=1}^M\sum_{\gamma\in [\Gamma]: j\in \supp(\mc K_{\gamma}^{(v)})}\|\mc K_\gamma^{(v)}\|\le g,
\end{align*}
which is called \emph{$g$-extensiveness} of $\mc K:= \sum_{v=1}^M \mc K_v$.

A wide class of physically relevant models, specifically Lindbladians generated by sparse Pauli strings, naturally fits into this framework. In the following lemma, we translate the Hamiltonian and jump operator parameters into the locality and extensiveness of the resulting superoperators.

\begin{lemma}
\label{lem:lindblad-mapping}
    Consider a Lindbladian $\mathcal{L} = \mathcal{H} + \sum_{v=1}^M \mathcal{D}[D_v]$ acting on $N$ qubits, where
    \begin{enumerate}
        \item the Hamiltonian $H = \sum_{\gamma} h_\gamma P_\gamma$ consists of Pauli strings with weight at most $k$; and 
        \item each jump operator $D_v = \sum_{\eta} d_{v,\eta} Q_{v,\eta}$ consists of Pauli strings with weight at most $k$.
    \end{enumerate}
    Then the total Lindbladian $\mathcal{L}$ can be decomposed as a sum of \emph{$2k$-local superoperators}:
    \begin{align*}
        \mathcal{L} = \sum_{\mu} \mathcal{K}_\mu, \quad \text{with } |\supp(\mathcal{K}_\mu)| \le 2k.
    \end{align*}
    Furthermore, let the Hamiltonian extensiveness be $g_H = \max_j \sum_{\gamma: j \in \supp(P_\gamma)} |h_\gamma|$ and the dissipative extensiveness be $g_D = \max_j \sum_{v} \sum_{\eta, \lambda: j \in \supp(Q_{v,\eta}) \cup \supp(Q_{v,\lambda})} |d_{v,\eta} d_{v,\lambda}^*|$.
    The system satisfies the \emph{$G$-extensiveness} condition with parameter
    \begin{align*}
        G = 2 g_H + 2 g_D.
    \end{align*}
\end{lemma}
\begin{proof}
    The Hamiltonian part $\mathcal{H} = -i \sum_\gamma h_\gamma [P_\gamma, \cdot]$ is a sum of terms with locality $|\supp(P_\gamma)| \le k$. Since $\|[P_\gamma, \cdot]\| \le 2\|P_\gamma\| = 2$, each term contributes at most $2|h_\gamma|$ to the local extensiveness.
    
    For the dissipative part, we expand $\mathcal{D}[D_v](\rho) = \sum_{\eta,\lambda} d_{v,\eta}d_{v,\lambda}^* \mathcal{S}_{\eta,\lambda}(\rho)$, where the basis superoperator is $\mathcal{S}_{\eta,\lambda}(\rho) = Q_{v,\eta}\rho Q_{v,\lambda}^\dagger - \frac{1}{2}\{Q_{v,\lambda}^\dagger Q_{v,\eta}, \rho\}$. 
    Each $\mathcal{S}_{\eta,\lambda}$ acts non-trivially only on $\supp(Q_{v,\eta}) \cup \supp(Q_{v,\lambda})$, so its locality is bounded by $2k$.
    Regarding the norm, since $Q_{v,\eta}$ and $Q_{v,\lambda}$ are Pauli strings with unit operator norm, standard norm inequalities for unitary conjugation and anti-commutators imply $\|\mathcal{S}_{\eta,\lambda}\|\le 1 + \frac{1}{2}(2) = 2$.
    Summing the coefficients weighted by this norm bound directly yields the extensiveness parameter $G = 2g_H + 2g_D$.
\end{proof}

This lemma ensures that the results we derive for general local superoperators are directly applicable to sparse Pauli-Lindbladians by simply substituting the locality $k \to 2k$ and extensiveness $g \to G$.

For general $k$-local Hamiltonians, the standard commutator bound defined in Eq.~\eq{def-comm} is $\alpha_{\rm comm}^{(q)} = O( q! (kg)^q)$ \cite{childs2021theory,mizuta2025commutator}, which causes the BCH series to diverge \cite{mizuta2025commutator}. This scaling of $\alpha_{\rm comm}^{(q)}$ remains the same for $k$-local superoperators, as the key requirement of the proof is locality rather than Hermiticity. \citet{mizuta2025commutator} provided an upper bound on the sum of nested commutators for a single Hamiltonian. However, the proof relies essentially on the commutativity structure determined by the supports of the operators, rather than the specific forms. Therefore, this result admits a straightforward generalization to a sequence of distinct local superoperators.

\begin{lemma}[Generalization of {\cite[Lemma~7]{mizuta2025commutator}}]
\label{lem:commu-bound}
    Let $\mc K_1, \ldots, \mc K_q$ be a sequence of superoperators on a size-$N$ lattice $\Lambda$ such that $\mc K_j = \sum_{v=1}^{M} \mc K_{j,v}$ is $k$-local and $g_j$-extensive for all $j\in [q  ]$. Then the following inequality on the right-nested commutators holds: \begin{align*}
        \sum_{v_1, \ldots, v_q=1}^{M} \big\|[\mc K_{1,v_1}, \ldots, \mc K_{q',v_{q'}}, \mc A, \mc K_{q'+1,v_{q'+1}}, \ldots, \mc K_{q,v_q}]\big\|\le q!(2k)^q\Big(\prod_{j=1}^qg_j\Big)\|\mc A\|,
    \end{align*}
    for any $q'\in[q]$.
\end{lemma}
A useful corollary can be obtained by replacing $\mc A$ in \lem{commu-bound} with a term from a $k$-local superoperator.
\begin{corollary}
\label{cor:commu-bound-sum}
    Let $\mc K_1, \ldots, \mc K_q$ be a sequence of superoperators on a size-$N$ lattice $\Lambda$ such that $\mc K_j = \sum_{v=1}^{M} \mc K_{j,v}$ is $k$-local and $g_j$-extensive for all $j\in [q  ]$. Then we have 
    \begin{align*}
        \sum_{v_1, \ldots, v_q=1}^{M} \big\|[\mc K_{1,v_1}, \ldots, \mc K_{q,v_q}]\big\|  \le N(q-1)!(2k)^{q-1}\prod_{j=1}^q g_j,
    \end{align*}
\end{corollary}
\begin{proof}
We apply \lem{commu-bound} to the superoperator sequence of length $q-1$, $\mc K_1, \ldots, \mc K_{q-1}$, and set the fixed operator $\mc A = \mc K_{q,v_q}$. Summing over the index $v_q$, we obtain
\begin{align*}
\sum_{v_1, \ldots, v_q=1}^{M} \big\|[\mc K_{1,v_1}, \ldots, \mc K_{q,v_q}]\big\| & \le \sum_{v_q=1}^{M} (q-1)! (2k)^{q-1} \Big(\prod_{j=1}^{q-1} g_j\Big) \|\mc K_{q,v_q}\| \\
&= (q-1)! (2k)^{q-1} \Big(\prod_{j=1}^{q-1} g_j\Big) \sum_{v_q=1}^{M}\|\mc K_{q,v_q}\|\\
&\le (q-1)!\Big(\prod_{j=1}^{q-1} 2k g_j\Big)  Ng_q \\
&=N(q-1)!(2k)^{q-1}\prod_{j=1}^q g_j,
\end{align*}
where the last inequality uses the property that $\mc K_q$ is $g_q$-extensive.
\end{proof}

We can now bound the doubly right-nested commutators defined in Eq.~\eq{def-beta-bar-general} using \lem{commu-bound}. 
\begin{theorem}\label{thm:doubly-nested-bound}
    Let $\mathcal{K} = \sum_{v=1}^{M} \mathcal{K}_{v}$ be a $k$-local and $g$-extensive superoperator on a lattice $\Lambda$, and let $\mathcal{A}$ be a superoperator with $|\supp(\mc A)|\le k$. For a given sequence of positive integers $q_1, q_2, \ldots, q_d$, let $P_{r} = 1+\sum_{j=r}^{d} q_j$ for $r=1, 2, \ldots, d$ and $P_{d+1}=1$. Then the following inequality holds for any position index $q' \in [q_d]$
    \begin{align}
    \label{eq:double-comm-k-local}
        \begin{multlined}
            \sum_{v_{1,1}, \ldots, v_{d, q_d}=1}^{M} \bigg\| \Big[ [\mathcal{K}_{v_{1,1}}, \ldots, \mathcal{K}_{v_{1, q_1}}], \ldots, [\mathcal{K}_{v_{d-1,1}}, \ldots, \mathcal{K}_{v_{d-1, q_{d-1}}}], \\
        [\mathcal{K}_{v_{d,1}}, \ldots, \mathcal{K}_{v_{d, q'}}, \mathcal{A}, \mathcal{K}_{v_{d, q'+1}}, \ldots, \mathcal{K}_{v_{d, q_d}}] \Big] \bigg\| 
        \end{multlined}\le \bigg(\prod_{r=1}^{d}P_{r+1} q_r!(2kg)^{q_r}\bigg)\|\mc A\|.
    \end{align}
\end{theorem}

\begin{proof}
    We first decompose each generator $\mathcal{K}_{v}$ into a sum of $k$-local terms. Specifically, let
    \begin{align*}
        \mc K_v = \sum_{\gamma=1}^{\Gamma} \mc K_{\gamma}^{(v)},
    \end{align*}
    where each $\mc K_{\gamma}^{(v)}$ acts non-trivially on at most $k$ sites. To simplify the sum, we re-index these terms as a single sequence $\hat{\mathcal{K}}_{j}$ for $j=1, \ldots, M\Gamma$, such that
    \begin{align*}
        \mc K = \sum_{v=1}^M \mc K_v = \sum_{j=1}^{M\Gamma} \hat{\mc K}_{j}.
    \end{align*}
    By the linearity of the commutator and the triangle inequality, the left-hand side of Eq.~\eqref{eq:double-comm-k-local} is bounded by the sum of the norms of the commutators of the constituent $k$-local terms
    \begin{align}
        \label{eq:triangle-ineq-step}
        \begin{multlined}
             \sum_{v_{1,1}, \ldots, v_{d, q_d}=1}^{M} \bigg\| \Big[ [\mathcal{K}_{v_{1,1}}, \ldots, \mathcal{K}_{v_{1, q_1}}], \ldots, 
        [\mathcal{K}_{v_{d,1}}, \ldots, \mathcal{K}_{v_{d, q'}}, \mathcal{A}_X, \ldots, \mathcal{K}_{v_{d, q_d}}] \Big] \bigg\| \\
        \le \sum_{u_{1,1}, \ldots, u_{d, q_d}=1}^{M\Gamma} \bigg\| \Big[ [\hat{\mathcal{K}}_{u_{1,1}}, \ldots, \hat{\mathcal{K}}_{u_{1, q_1}}], \ldots, 
        [\hat{\mathcal{K}}_{u_{d,1}}, \ldots, \hat{\mathcal{K}}_{u_{d, q'}}, \mathcal{A}, \ldots, \hat{\mathcal{K}}_{u_{d, q_d}}] \Big] \bigg\|.
        \end{multlined}
    \end{align}
    For any layer $r \in [d]$, let $\vec{u}_r = (u_{r,1}, \ldots, u_{r, q_r})$ denote the vector of indices ranging from $1$ to $M\Gamma$. We define the nested commutator for layer $r$ as
    \begin{align*}
        \hat{\mathcal{K}}_{\vec{u}_r} := [\hat{\mathcal{K}}_{u_{r,1}}, \ldots, \hat{\mathcal{K}}_{u_{r, q_r}}].
    \end{align*}
    Similarly, let $\hat{\mathcal{W}}_{\vec{u}_d} := [\hat{\mathcal{K}}_{u_{d,1}}, \ldots, \mathcal{A}_X, \ldots, \hat{\mathcal{K}}_{u_{d, q_d}}]$ denote the innermost term containing $\mc A$ at the $q'$-th position.
    The RHS of Eq.~\eqref{eq:triangle-ineq-step} can be rewritten as
    \begin{align*}
        \sum_{\vec{u}_{1}, \ldots, \vec{u}_{d}} \bigg\| \Big[ \hat{\mathcal{K}}_{\vec{u}_1}, \ldots, \hat{\mathcal{K}}_{\vec{u}_{d-1}}, 
        \hat{\mathcal{W}}_{\vec{u}_d} \Big] \bigg\|.
    \end{align*}
    Define the partial nested commutator starting from layer $r$ as
    \begin{align*}
        \mathcal{W}_{\vec{u}_{r}, \ldots, \vec{u}_d} := \Big[ \hat{\mathcal{K}}_{\vec{u}_r}, \ldots, \hat{\mathcal{K}}_{\vec{u}_{d-1}}, \hat{\mathcal{W}}_{\vec{u}_d} \Big].
    \end{align*}
    We prove the following bound by backward induction on $r$ from $d$ to $1$
    \begin{align} \label{eq:induct-W}
        \sum_{\vec{u}_{r}, \ldots, \vec{u}_{d}} \Big\|\mathcal{W}_{\vec{u}_{r}, \ldots, \vec{u}_d}\Big\|\le \bigg(\prod_{j=r}^{d}P_{j+1} q_j!(2kg)^{q_j}\bigg)\|\mc A\|.
    \end{align}
    
    Consider the base case $r=d$. The term is $\sum_{\vec{u}_d} \|\hat{\mathcal{W}}_{\vec{u}_d}\|$. Since $\mathcal{W}_{\vec{u}_d}$ involves one fixed operator $\mathcal{A}$ and $q_d$ operators from a $k$-local set, we apply \lem{commu-bound}. This yields the bound $q_d!(2kg)^{q_d}\|\mathcal{A}\|$. Since $P_{d+1}=1$, the base case holds.

    Assume Eq.~\eqref{eq:induct-W} holds for $r$. We consider the sum for $r-1$
    \begin{align*}
        \sum_{\vec{u}_{r-1}, \ldots, \vec{u}_d} \Big\|\mathcal{W}_{\vec{u}_{r-1}, \ldots, \vec{u}_d}\Big\| 
        &= \sum_{\vec{u}_{r}, \ldots, \vec{u}_d} \sum_{\vec{u}_{r-1}} \bigg\| \Big[ \hat{\mathcal{K}}_{\vec{u}_{r-1}}, \mathcal{W}_{\vec{u}_{r}, \ldots, \vec{u}_d} \Big] \bigg\|.
    \end{align*}
    Let $X_{\mathcal{W}} = \supp(\mathcal{W}_{\vec{u}_{r}, \ldots, \vec{u}_d})$. Its size is bounded by $|X_{\mathcal{W}}| \le k P_r$. The inner commutator is nonzero only if at least one operator in $\hat{\mathcal{K}}_{\vec{u}_{r-1}}$ overlaps with $X_{\mathcal{W}}$.
    We explicitly expand the sum over $\vec{u}_{r-1}$ by iterating over which operator overlaps with $X_{\mathcal{W}}$. Let $j \in [q_{r-1}]$ be the index of this operator. We can upper bound the sum by
    \begin{align*}
        & \sum_{\vec{u}_{r-1}} \bigg\| \Big[ \hat{\mathcal{K}}_{\vec{u}_{r-1}}, \mathcal{W}_{\vec{u}_{r}, \ldots, \vec{u}_d} \Big] \bigg\| \\
        \le~& \sum_{j=1}^{q_{r-1}} \sum_{\substack{u_{r-1, j} \in [M\Gamma] \\ \supp(\hat{\mathcal{K}}_{u_{r-1, j}}) \cap X_{\mathcal{W}} \neq \emptyset}} \sum_{\substack{u_{r-1, k} \in [M\Gamma] \\ (\forall k \neq j)}} \Big\| \big[ \hat{\mathcal{K}}_{u_{r-1,1}}, \ldots, \hat{\mathcal{K}}_{u_{r-1, q_{r-1}}}, \mathcal{W}_{\vec{u}_{r}, \ldots, \vec{u}_d} \big] \Big\| \\
        \le~& \sum_{j=1}^{q_{r-1}} \sum_{\substack{u_{r-1, j} \in [M\Gamma] \\ \supp(\hat{\mathcal{K}}_{u_{r-1, j}}) \cap X_{\mathcal{W}} \neq \emptyset}} (q_{r-1}-1)! (2kg)^{q_{r-1}-1} \|\hat{\mathcal{K}}_{u_{r-1, j}}\| \cdot 2 \Big\|\mathcal{W}_{\vec{u}_{r}, \ldots, \vec{u}_d}\Big\| \\
        \le~& \sum_{j=1}^{q_{r-1}} (k P_r g) \cdot (q_{r-1}-1)! (2kg)^{q_{r-1}-1} \cdot 2 \Big\|\mathcal{W}_{\vec{u}_{r}, \ldots, \vec{u}_d}\Big\| \\
        =~& q_{r-1} \cdot (k P_r g) \cdot (q_{r-1}-1)! (2kg)^{q_{r-1}-1} \cdot 2 \Big\|\mathcal{W}_{\vec{u}_{r}, \ldots, \vec{u}_d}\Big\| \\
        =~& P_r \cdot q_{r-1}! \cdot (2kg)^{q_{r-1}} \Big\|\mathcal{W}_{\vec{u}_{r}, \ldots, \vec{u}_d}\Big\|.
    \end{align*}
    Here, the second inequality follows from applying \lem{commu-bound} to the sum over the $q_{r-1}-1$ operators not at index $j$ with $\mc A:= \hat{\mathcal{K}}_{u_{r-1, j}}$.  The third inequality uses the $g$-extensiveness property: summing the norm of $\hat{\mathcal{K}}_{u_{r-1, j}}$ over all operators overlapping with $X_{\mathcal{W}}$ yields at most $|X_{\mathcal{W}}|g \le k P_r g$.
    
    Substituting this back and using the inductive hypothesis leads to
    \begin{align*}
        \sum_{\vec{u}_{r-1}, \ldots, \vec{u}_d} \Big\|\mathcal{W}_{\vec{u}_{r-1}, \ldots, \vec{u}_d}\Big\| 
        &\le P_r q_{r-1}! (2kg)^{q_{r-1}} \sum_{\vec{u}_{r}, \ldots, \vec{u}_d} \Big\|\mathcal{W}_{\vec{u}_{r}, \ldots, \vec{u}_d}\Big\| \\
        &\le P_r q_{r-1}! (2kg)^{q_{r-1}}  \bigg(\prod_{j=r}^{d}P_{j+1} q_j!(2kg)^{q_j}\bigg)\|\mc A\|.
    \end{align*}
    This completes the induction.
\end{proof}

Plugging this bound into \thm{general-bch-truncation}, we can bound the truncation error for a $k$-local system as follows. Note that if all $\mathcal{K}_v$ are Hamiltonians, this bound matches the one in \citet[Theorem 3]{mizuta2025commutator}. However, we use a distinct proof which applies to general superoperators beyond lattice systems.
\begin{lemma}
\label{lem:bch-approx} 
    For any $\eps\in (0,1/2)$ and  $q_0 \ge \log(2N/\eps)$, if $4e^2kgq_0\le 1$, then we have \begin{align*}
        \left\|\exp\Big(\sum_{q=1}^{q_0} \Phi_q\Big) - \prod_{v = 1}^Me^{\mc K_v} \right\|\le \eps.
    \end{align*}
\end{lemma}
\begin{proof} 
    For any positive integer sequence $q_1, \ldots, q_{d+1}$, let 
    \begin{align*}
        P_{i} = \sum_{j=i}^{d+1} q_j,
    \end{align*}
    for $i\le d+1$, $P_{d+2} = 1$, and $q = q_1+\cdots+q_{d+1}$. Then, we have 
    \begin{align*}
    &\bar{\beta}_{\comm}^{(q_1, \ldots, q_{d+1})}\\
    =~& \sum_{v_{d+1, q_{d+1}}=1}^{M} \sum_{v_{1,1}, \ldots, v_{d+1, q_{d+1}-1} = 1}^M \bigg\| \Big[ 
         [\mc K_{v_{1,1}}, \ldots, \mc K_{v_{1,q_1}}], \ldots, 
        [\mc K_{v_{d+1,1}}, \ldots, \mc K_{v_{d+1,q_{d+1}}}]
     \Big] \bigg\| \\
    \le~& \bigg(\prod_{i=1}^{d+1}P_{i+1}q_i!(2kg)^{q_i }\bigg)\frac{1}{g}\sum_{v_{d+1, q_{d+1}}=1}^{M} \|\mc K_{v_{d+1, q_{d+1}}}\|\\
    =~&  \bigg(\prod_{i=1}^{d+1}P_{i+1}q_i!(2kg)^{q_i }\bigg)N,
\end{align*}
where the third line follows from \thm{doubly-nested-bound}, and the fourth line follows from 
\begin{align*}
    \sum_{v_{d+1,q_{d+1}}=1}^{M}\|\mc K_{v_{d+1, q_{d+1}}}\|\le \sum_{x\in\Lambda}\sum_{\substack{v = 1, \ldots, M \\ \supp(\mc K_v)\ni x}} \|\mc K_{v}\|\le gN.
\end{align*}
For $q_1, \ldots, q_{d+1}\le q_0$, we have 
\begin{align*}
    \bar{\beta}_{\comm}^{(q_1, \ldots, q_{d+1})}\le q^{d}\bigg(\prod_{i=1}^{d+1}q_0^{q_{i}}(2kg)^{q_i }\bigg)N = q^{d} (2q_0kg)^q N,
\end{align*}
where the inequality follows from $P_{i+1}\le q$ (noting that $P_{d+2} = 1$ does not contribute a factor of $q$). 
Then, we have 
\begin{align*}
    \bar{\beta}_{\comm,q_0} &\le (2q_0 kg)^{q_0+1} N+ \sum_{r = 0}^{\infty}\frac{q^{r} }{(r+1)!} \sum_{q=q_0+2}^{\infty}\sum_{\substack{1\le p_1, \ldots, p_{r+1}\le q_0\\p_1+\cdots+p_{r+1} = q}}(2q_0 kg)^q N \\ 
    &\le  (2q_0 kg)^{q_0+1} N+  \sum_{q=q_0+2}^{\infty}\sum_{\substack{1\le p_1, \ldots, p_{r+1}\le q_0\\p_1+\cdots+p_{r+1} = q}} (2eq_0 kg)^q N\\
    &\le(2q_0 kg)^{q_0+1} N+  \sum_{q=q_0+2}^{\infty} (4eq_0 kg)^q N   \\
    &\le (2q_0kg)^{q_0+1} N + (4eq_0 kg)^{q_0+1} N \sum_{q = 1}^{\infty} 2^{-q} \\
    &\le 2(4eq_0kg)^{q_0+1}N \le (4e^2q_0 kg)^{q_0+1}\le 1,
\end{align*}
where the third line follows from 
\begin{align*}
    \sum_{\substack{1\le q_1, \ldots, q_{r+1}\le q_0\\q_1+\cdots q_{r+1}=q}}1\le  \sum_{\substack{1\le q_1, \ldots, q_{r+1}\\q_1+\cdots q_{r+1}=q}}1=\binom{q-1}{r}\le 2^q,
\end{align*}
and the fourth line follows from $e^{-q_0-1} \le \eps/(2N) \le 1/(2N)$.
As $\bar{\beta}_{\comm,q_0}\le 1$ is satisfied, we can apply \thm{general-bch-truncation} to obtain 
\begin{align*}
\left\|\exp\Big(\sum_{q=1}^{q_0} \Phi_q\Big) - \prod_{v = 1}^Me^{\mc K_v} \right\|
\le  e\bar{\beta}_{\comm,q_0}
\le 2e(4eq_0kg)^{q_0+1}N\le 2e^{-q_0}N\le \eps.
\end{align*}
\end{proof}
Based on the convergence condition in Eq.~\eq{converg-bch} and the bound of $\alpha_{\rm comm}$ in Eq.~\eq{trivial-alpha}, the BCH series $\sum_{q=1}^{\infty} \Phi_q$ is guaranteed to converge only when $\sum_{v=1}^M \|\mc K_v\| = O(1)$, which requires $g = O(1/N)$. In contrast, \lem{bch-approx}  shows that the exponential of the truncated BCH formula approximates the exponential product in a much larger regime of $g = O(1/(k\log N))$, where the BCH series is not guaranteed to converge. 

%%%%%%%%%%%%%%%%%%%%%%%%%%%%%%%%%%%%%%%%%%%%%%%%%%%%%%%%%%%%%%%%%%%%%%%%%%%%%%%%%%%%%%%%%%%%%%%%%%%%%%%%%%%%%%%%

\subsection{Application to second-quantized electronic structure}
Following the notation of \citet{childs2021theory}, we consider the second-quantized electronic structure Hamiltonian in the plane-wave dual basis. In this representation, the Hamiltonian for a system with $n$ spin-orbitals takes the form
\begin{align}
    \label{eq:def-tuv}
    H = \underbrace{\sum_{j,k=1}^n t_{j,k} A_j^\dagger A_k}_{T} 
      + \underbrace{\sum_{j=1}^n u_j N_j}_{U} 
      + \underbrace{\sum_{j,k=1}^n g_{j,k} N_j N_k}_{V},
\end{align}
where $A_j^\dagger, A_j$ are the creation and annihilation operators satisfying $\{A_j, A_k^\dagger\} = \delta_{jk}$, and $N_j = A_j^\dagger A_j$. Here, $T$ represents the kinetic energy with hopping amplitudes $t_{j,k}$, $U$ denotes the single-particle potential with coefficients $u_j$, and $V$ describes the two-body interaction with strength $g_{j,k}$.
Note that $\|A_j^{\dagger}\|=\|A_k\|=\|N_j\|=1$. We can apply the triangle inequality to upper bound the norms $\|T\|$, $\|U\|$, and $\|V\|$ by the $1$-norm of the corresponding coefficient vectors $\|\vec{t}\|_1$, $\|\vec{u}\|_1$, and $\|\vec{v}\|_1$. \citet{childs2021theory} proved the follow upper bounds for the $1$-norm and $\infty$-norm of the coefficient vectors. 
\begin{lemma}[{\cite[Proposition F.2]{childs2021theory}}]
\label{lem:norm-tuv}
    The $1$-norm and $\infty$-norm of the coefficients vectors of the kinetic operator $T$ and potential operators $U$, $V$ in Eq.~\eq{def-tuv} can be bounded as \begin{align*}
        &\|\vec{t}\|_{\infty} = O(1/n), && \|\vec{t}\|_1 = O(n) \\
        &\|\vec{u}\|_{\infty} = O(n), && \|\vec{u}\|_1 = O(n^2) \\
        &\|\vec{g}\|_{\infty} = O(1), && \|\vec{g}\|_1  = O(n^2). 
    \end{align*}
\end{lemma}
Denote the summand in Eq.~\eq{def-tuv} by $H_{v}$, which gives \begin{align*}
    H = \sum_{v = 1}^{M} H_{v},
\end{align*}
with $M = 2n^2 +n$.
The right-nested commutators of $H_v$ can be expressed as \begin{align*}
    W = \sum_{\vec{j}, \vec{k},\vec{\ell}}w_{\vec{j}, \vec{k},\vec{\ell}}\cdots (A_{j_x}^{\dagger} A_{k_x}) \cdots (N_{\ell_y})\cdots,
\end{align*}
where $\vec{j}, \vec{k},\vec{\ell}$ are vectors of orbitals. The largest number of $A_{j_x}^{\dagger} A_{k_x}$ and $N_{\ell_y}$ terms in the summand is called the ``layer'' of $W$. For any $W$, the norms of the commutators of terms in $T$, $U$, and $V$ with $W$ are bounded by the $\infty$-norms of $t$, $u$, and $g$, and the $1$-norm of $w$.
\begin{lemma}[{\cite[Proposition F.4]{childs2021theory}}]
\label{lem:recursion-W}
     Let an electronic-structure Hamiltonian be given as in Eq.~\eq{def-tuv}. The following statements hold for a general second-quantized operator $W$ with $q$ layers:
     \begin{enumerate}
         \item $\tilde{W}=[T, W]$ is an operator with $q$ layers and $\|\vec{\tilde{w}}\|_1 \leq 2 q n\|\vec{t}\|_{\infty}\|\vec{w}\|_1 $;
         \item $\tilde{W}=[U, W]$ is an operator with $q$ layers and $\|\vec{\tilde{w}}\|_1 \leq 2 q\|\vec{u}\|_{\infty}\|\vec{w}\|_1$; and 
         \item $\tilde{W}=[V, W]$ is an operator with $q+1$ layers and $\|\vec{\tilde{w}}\|_1 \leq 4 q n\|\vec{g}\|_{\infty}\|\vec{w}\|_1$.
     \end{enumerate}
\end{lemma}
A Corollary of this lemma is an upper bound for the sum of right-nested commutator norms. 
\begin{corollary}
\label{cor:right-comm-second}
    The sum of the norms of the right-nested commutators of $H_v$ satisfies 
    \begin{align*}
        \sum_{v_1, \ldots, v_q=1}^{M} \|[H_{v_1}, \ldots, H_{v_q}]\| = \mathcal{O}(q! n^q),
    \end{align*}
    for any $q\ge 2$. In addition, each $[H_{v_1}, \ldots, H_{v_q}]$ has at most $q$ layers. 
\end{corollary}
\begin{proof}
    We prove the statement by induction on $q$. 
    
    Consider the base case $q=2$. The LHS is $\sum_{v_1, v_2=1}^{M} \|[H_{v_1}, H_{v_2}]\|$, where each $H_{v_1}$ and $H_{v_2}$ is a term belonging to $T$, $U$, or $V$. Note that terms from $U$ and $V$ commute with each other, so we only need to consider commutators between $T$ and $\{U,V\}$. This yields
    \begin{align*}
        \sum_{v_1, v_2=1}^{M} \|[H_{v_1}, H_{v_2}]\| 
        &= \mathcal{O}\big(n\|\vec{t}\|_{\infty} \|\vec{u}\|_1+ n\|\vec{t}\|_{\infty} \|\vec{g}\|_1 + n\|\vec{t}\|_{\infty} \|\vec{t}\|_1\big) \\
        &= \mathcal{O}(n^2).
    \end{align*}
    The first equality follows from \lem{recursion-W} and the second follows from the coefficient norms in \lem{norm-tuv}. In addition, the layer of $[H_{v_1}, H_{v_2}]$ is at most $2$. Therefore, the base case holds. 
    
    Assume the statement holds for $q-1$. Consider the sum for $q$: 
    \begin{align*}
         \sum_{v_1, \ldots, v_q=1}^{M} \|[H_{v_1}, \ldots, H_{v_q}]\| = \sum_{v_1=1}^M  \sum_{v_2, \ldots, v_q=1}^{M} \Big\| \big[H_{v_1}, [H_{v_2}, \ldots, H_{v_q}] \big] \Big\|. 
    \end{align*} 
    By the induction hypothesis, $[H_{v_2}, \ldots, H_{v_q}]$ has at most $q-1$ layers.
    According to \lem{recursion-W}, commuting with terms belonging to $T$, $U$, or $V$ increases the norm by a factor proportional to $(q-1)n\|\vec{t}\|_\infty = \mathcal{O}(q)$, $(q-1)\|\vec{u}\|_\infty =\mathcal{O}(qn)$, or $(q-1)n\|\vec{g}\|_\infty=\mathcal{O}(qn)$, respectively. Summing over all possible $H_{v_1}$, we have
    \begin{align*}
        \sum_{v_1=1}^M \sum_{v_2, \ldots, v_q=1}^{M} \Big\| \big[H_{v_1}, [H_{v_2}, \ldots, H_{v_q}] \big] \Big\| 
        = \mathcal{O}\Big(qn \sum_{v_2, \ldots, v_q=1}^{M} \|[H_{v_2}, \ldots, H_{v_q}]\|\Big) 
        = \mathcal{O}(q!n^{q}),
    \end{align*}
    where the last equality follows from the induction hypothesis. By \lem{recursion-W}, commuting with an $H_v$ increases the layer by at most one, and hence $[H_{v_1}, \ldots, H_{v_q}]$ has at most $q$ layers.
\end{proof}



Let $\mc{K}_v:=-\i \ad_{H_v}$ and \begin{align}
\label{eq:def-K-second-quantized}
    \mc{K}:= \sum_{v=1}^M  \mc K_v= -\i \ad_H. 
\end{align}
We now apply \thm{general-bch-truncation} to bound the truncation error of the BCH formula of $\prod_{v=1}^M\mc K_v$.
\begin{lemma}
    Let $\mc K = \sum_{v=1}^M \mc K_v$ be the Liouvillian in Eq.~\eq{def-K-second-quantized}. For a given sequence of positive integers $q_1, q_2, \ldots, q_d$, let $P_{r} = \sum_{j=r}^{d} q_j$ for $r=1, 2, \ldots, d$ and $P_{d+1}=1$. Then, the doubly right-nested commutator bound $\bar{\beta}_{\comm}^{(q_1, \ldots, q_{d})}$ for $\mc K$ (see Eq.~\eq{def-beta-bar-general}) satisfies
    \begin{align}
    \label{eq:double-comm-second-quantized}
        \bar{\beta}_{\comm}^{(q_1, \ldots, q_{d})}= O\bigg(\prod_{r=1}^{d}P_{r+1} q_r!n^{q_r}\bigg).
    \end{align}
\end{lemma}
\begin{proof}
    First, we reduce the problem to bounding commutators of Hamiltonian terms. By applying the Jacobi identity $[\ad_{X}, \ad_Y] = \ad_{[X,Y]}$ recursively and using $\|\ad_X\| \le 2\|X\|$, the doubly right-nested commutator bound of $\mc K_v$ satisfies
    \begin{align*}
    \bar{\beta}_{\comm}^{(q_1, \ldots, q_{d})}
    &= \sum_{v_{1,1}, \ldots, v_{d, q_{d}} = 1}^M \bigg\| \ad_{\Big[ 
         [H_{v_{1,1}}, \ldots, H_{v_{1,q_1}}], \ldots, 
        [H_{v_{d,1}}, \ldots, H_{v_{d,q_{d}}}]
     \Big]} \bigg\|\\
     &\le 2\sum_{v_{1,1}, \ldots, v_{d, q_{d}} = 1}^M \bigg\|{\Big[ 
         [H_{v_{1,1}}, \ldots, H_{v_{1,q_1}}], \ldots, 
        [H_{v_{d,1}}, \ldots, H_{v_{d,q_{d}}}]
     \Big]} \bigg\|.
    \end{align*}
    Therefore, it suffices to bound the sum of doubly right-nested commutators of $H_v$. 
    
    To simplify the notation, for any layer $r \in [d]$, let $\vec{v}_r = (v_{r,1}, \ldots, v_{r, q_r})$ denote the index vector, and define the nested commutator for that layer as
    \begin{align*}
        H_{\vec{v}_r} := [H_{v_{r,1}}, \ldots, H_{v_{r, q_r}}].
    \end{align*}
    We define the partial nested commutator starting from layer $r$ as
    \begin{align*}
        W_{\vec{v}_{r}, \ldots, \vec{v}_d} := \Big[ H_{\vec{v}_r}, \ldots, H_{\vec{v}_{d-1}}, H_{\vec{v}_d} \Big].
    \end{align*}
    We prove the following bound by backward induction on $r$ from $d$ to $1$:
    \begin{align} \label{eq:induct-W-elec}
        \sum_{\vec{v}_{r}, \ldots, \vec{v}_{d}} \Big\|W_{\vec{v}_{r}, \ldots, \vec{v}_d}\Big\| = \mathcal{O}\bigg(\prod_{j=r}^{d}P_{j+1} q_j! n^{q_j}\bigg).
    \end{align} In addition, we show that $W_{\vec{v}_{r}, \ldots, \vec{v}_d}$ has at most $P_r$ layers.  
    
    Consider the base case $r=d$. By \cor{right-comm-second}, the LHS of Eq.~\eq{induct-W-elec} is \begin{align*}
        \sum_{\vec{v}_d} \|H_{\vec{v}_d}\| = \mc{O}(q_d ! n^{q_d}) = \mc{O}(P_{d+1}q_d ! n^{q_d}),
    \end{align*} and $H_{\vec{v}_d}$ has $q_d$ layers, so the base case holds. 

    Assume Eq.~\eqref{eq:induct-W-elec} holds for $r$. We consider the sum for $r-1$:
    \begin{align*}
        \sum_{\vec{v}_{r-1}, \ldots, \vec{v}_d} \Big\|W_{\vec{v}_{r-1}, \ldots, \vec{v}_d}\Big\| 
        &= \sum_{\vec{v}_{r}, \ldots, \vec{v}_d} \sum_{\vec{v}_{r-1}} \bigg\| \Big[ H_{\vec{v}_{r-1}}, W_{\vec{v}_{r}, \ldots, \vec{v}_d} \Big] \bigg\|.
    \end{align*}
    If $q_{r-1}\ge 2$, we have \begin{align*}
        \sum_{\vec{v}_{r}, \ldots, \vec{v}_d} \sum_{\vec{v}_{r-1}} \bigg\| \Big[ H_{\vec{v}_{r-1}}, W_{\vec{v}_{r}, \ldots, \vec{v}_d} \Big] \bigg\| &\le 2\Big(\sum_{\vec{v}_{r-1}}\|H_{\vec{v}_{r-1}}\|\Big)\Big(\sum_{\vec{v}_{r}, \ldots, \vec{v}_d}\|W_{\vec{v}_{r}, \ldots, \vec{v}_d}\|\Big) \\
        &= \mc{O}\Big(q_{r-1}!n^{q_{r-1}}\prod_{j=r}^{d}P_{j+1} q_j! n^{q_j}\Big)\\
        &= \mc{O}\Big(\prod_{j=r-1}^{d}P_{j+1} q_j! n^{q_j}\Big),
    \end{align*}
    where the second line follows from \cor{right-comm-second} and the induction hypothesis. The layer count of $W_{\vec{v}_{r-1}, \ldots, \vec{v}_d}$ is the sum of the layer counts of $H_{\vec{v}_{r-1}}$ and $W_{\vec{v}_{r}, \ldots, \vec{v}_d}$, which is $q_{r-1}+ P_{r} = P_{r-1}$. 
    
    Then we consider the case of $q_{r-1}=1$. We have 
    \begin{align*}
        \sum_{\vec{v}_{r-1}, \ldots, \vec{v}_d} \Big\|W_{\vec{v}_{r-1}, \ldots, \vec{v}_d}\Big\| 
       = \sum_{v}\sum_{\vec{v}_{r}, \ldots, \vec{v}_d} \Big\|[H_{v},W_{\vec{v}_{r}, \ldots, \vec{v}_d}]\Big\|.
    \end{align*}
    By the induction hypothesis, $W_{\vec{v}_{r}, \ldots, \vec{v}_d}$ has at most $P_{r}$ layers. Using the same argument in the proof of \cor{right-comm-second}, commuting with an $H_v$ increases the norm by a factor of $\mc{O}(P_r n)$. Therefore, we have \begin{align*}
        \sum_{v}\sum_{\vec{v}_{r}, \ldots, \vec{v}_d} \Big\|[H_{v},W_{\vec{v}_{r}, \ldots, \vec{v}_d}]\Big\| &= \mc{O}\Big(P_r n \sum_{\vec{v}_{r}, \ldots, \vec{v}_d} \|W_{\vec{v}_{r}, \ldots, \vec{v}_d}\|\Big)\\
        &= \mc{O}\Big(P_r n \prod_{j=r}^{d}P_{j+1} q_j! n^{q_j}\Big) \\
        &= \mc{O}\Big( \prod_{j=r-1}^{d}P_{j+1} q_j! n^{q_j}\Big)
    \end{align*}
    where the second line follows from the induction hypothesis. By \lem{recursion-W}, $W_{\vec{v}_{r-1}, \ldots, \vec{v}_d}$ has at most $1+P_{r} = P_{r-1}$ layers.  
    This completes the induction. Setting $r=1$ yields the lemma statement.
\end{proof}


\section{Analysis of the extrapolation error}
Given a Hamiltonian $H = \sum_{\gamma=1}^{\Gamma} H_{\gamma}$, an initial state $\rho_0$, and an observable $O$, we aim to estimate \begin{align*}
    \tr[Oe^{-\i \ad_H T} \rho_0],
\end{align*}
for any simulation time $T$.
Suppose that the Hamiltonian evolution driven by each $H_{\gamma}$ can be efficiently implemented by a layer of gates. Then, for sufficiently small $t$, we can implement the Hamiltonian evolution $e^{-\i H t}$ approximately by a $p$-th order product formula of the form
\begin{align*}
    P(t) = \prod_{v=1}^Me^{-\i t a_{v} H_{\gamma_v}},
\end{align*}
where the product is ordered from right to left. Here, $a_{v}$ are real coefficients, $M = \Upsilon\Gamma$ is the total number of unitary gates, and the sequence $\gamma_v$ specifies which Hamiltonian is applied at step $v$. Denote $a_{\max} = \max_{v\in [M]} |a_v|$. The corresponding ideal channel $\P(t)$ is a superoperator given by
\begin{align*}
    \P(t) = \prod_{v=1}^Me^{-\i t a_{v} \ad_{H_{\gamma_v}}},
\end{align*}
where $\A_{\gamma}:=\ad_{H_{\gamma}}$ denotes the adjoint action of $H_{\gamma}$. 

Suppose that the system is subject to physical noise with a controllable strength $\lambda$. Specifically, after each layer of unitary gates, a noise channel $e^{\lambda \mc L_{v}}$ is applied, where $\mc L_{v}$ is a Lindbladian. The resulting noisy channel is
\begin{align*}
    \tilde{\P}(t,\lambda)=\prod_{v=1}^M\left( e^{\lambda \mc L_{v}} e^{-\i t a_{v} \ad_{H_{\gamma_v}}} \right).
\end{align*} 
In experimental implementations, the noise strength is lower-bounded by the device's native noise level, denoted by $\lambda_{\min}$. \xz{add some overview on how to tune the noise level like unitary folding or learning and recovering the Pauli channel.} Therefore, we assume that $\lambda$ can be varied within the range $[\lambda_{\min}, \infty)$. 
Given $s\in (0,1)$ such that $1/s\in \mathbb{Z}$, we implement $\mc P(sT)^{1/s}$ with initial state $\rho_0$ and then measure the observable $O$ to approximate $\tr[Oe^{-\i \ad_H T}\rho_0]$. We denote the expectation of the measurement results with noisy circuit by  \begin{align*}
    g(s,\lambda) := \tr[O\tilde{\mc P}(sT,\lambda)^{1/s}\rho_0],
\end{align*}
and note that \begin{align*}
    \lim_{s,\lambda\to 0} g(s,\lambda) = \tr[Oe^{-\i \ad_H T}\rho_0]
\end{align*}
is the quantity we want. To obtain this limit value, one could perform two successive extrapolation: first, use ZNE to extrapolate $\lambda \to 0$ for a fixed $s$, and then use Trotter error mitigation \cite{watson2025exponentially} to extrapolate $s \to 0$ using the ZNE output. Although this method is possible, it involves extrapolating two independent parameters, which requires more sample points. Furthermore, each extrapolation step amplifies other sources of noise, such as state preparation and measurement (SPAM) errors, making the method less robust.

Therefore, we consider a more efficient and robust approach: a unified 1D extrapolation. By coupling the noise strength and the step size, we reduce this 2D problem to a single-variable extrapolation. In general, we could set $\lambda$ to be an analytic function of $s$, $\lambda(s)$, such that $\lambda(0)=0$. This ensures that the extrapolated function $g(s,\lambda(s))$ has the correct limit value at $s = 0$. Furthermore, its analyticity near $s=0$ guarantees the existence of a power series expansion, which is the fundamental requirement for Richardson extrapolation. For general analytic function $\lambda(s) =\sum_{j=p_0+1}^{\infty} c_j (sT)^j$, the noise error will be dominated by its leading-order term, $O(s^{p_0+1})$. Thus, for analyzing the asymptotic scaling, it is sufficient to consider a simple monomial form  $\lambda(s) = c(sT)^{p_0+1}$ for some $c>0$ and $p_0\in \mb{Z}_+$.

We now determine the optimal exponent $p_0$. The total error of the estimator $g(s, \lambda(s))$ arises from two dominant sources: the algorithmic (Trotter) error, which scales as $O(s^{p+1})$ for a $p$-th order product formula, and the physical noise error, which scales linearly with $\lambda$ in the leading order, and thus scales as $O(s^{p_0+1})$.
The leading-order error of the combined system is governed by the term with the slower decay, i.e., $O(s^{\min\{p+1, p_0+1\}})$. To maximize the efficiency of the extrapolation, we aim to maximize this leading-order exponent. If we choose $p_0 < p$, the noise error dominates ($O(s^{p_0+1})$) and becomes the bottleneck, wasting the precision of the $p$-th order Trotter formula. Conversely, if $p_0 > p$, the Trotter error dominates ($O(s^{p+1})$), implying that we are suppressing the noise strength more aggressively than necessary.
The optimal strategy is to match the scaling of the two error sources by setting $p_0 = p$. This choice ensures that both the leading-order Trotter error and the leading-order noise error appear at the same order, $O(s^{p+1})$, allowing the 1D Richardson extrapolation to eliminate both simultaneously and efficiently. We further remark that this matching strategy is not merely a heuristic. As we will demonstrate in our applications (\xz{add ref to application}), this choice enables our protocol to achieve a simulation depth that saturates the fundamental limits of quantum error mitigation under Pauli noise, as established in Refs.~\cite{wang2021noise, cai2021multi}.

To analyze the Richardson extrapolation error, our goal is to bound the remainder of the power series expansion of \begin{align*}
    f(s):= g(s, c(sT)^{p+1}) =  \tr[O\tilde{\mc P}(sT,c(sT)^{p+1})^{1/s}\rho_0].
\end{align*}
The $1/s$ exponent in the definition of $f(s)$ makes a direct power series analysis difficult. The standard approach is to first analyze the single-step propagator, $\tilde{\mc P}(sT,c(sT)^{p+1})$, by finding its effective generator. Note that the full channel can be written as \begin{align*}
    \tilde{\mc P}(t,\lambda) = \prod_{v=1}^{2M} e^{\mc K_{v}},
\end{align*} where each $\mc K_v$ is a Lindbladian defined as 
\begin{align}
 	\mathcal{K}_j= \begin{cases}-\i ta_{v}\ad_{H_{\gamma_v}} & \text{if } j = 2v-1\\ \lambda\mathcal{L}_v & \text{if } j = 2v \end{cases}, 	\label{eq:def-K}
\end{align}
for $v \in [M]$. Therefore, the effective generator of $\tilde{\mc P}(t,\lambda)$ is given by the BCH formula
\begin{align*}
    \tilde{\mc P}(t,\lambda)  = \exp\left(\sum_{q=1}^\infty \Phi_q(t,\lambda)\right),
\end{align*}
where $\Phi_q$ is the $q$-th order term in the BCH formula in Eq.~\eq{bch} with $2M$ generators $\mc K_1, \ldots, \mc K_{2M}$ defined in Eq.~\eq{def-K}.
The full evolution is then \begin{align*}
    \tilde{\mc P}(sT,c(sT)^{p+1})^{1/s} = \exp(\frac{1}{sT}\sum_{q=1}^\infty \Phi_q(sT,c(sT)^{p+1}) T). 
\end{align*}
The analysis of the power series expansion of $\tilde{\mc P}(sT,c(sT)^{p+1}))^{1/s}$ reduces to analyzing the scaled effective generator \begin{align*}
    \frac{1}{sT} \sum_{q=1}^\infty \Phi_q(sT,c(sT)^{p+1}).
\end{align*}
To resolve the divergence issue of this BCH series as discussed in \sec{bch-truncate}, we consider the truncated scaled effective generator \begin{align}
\label{eq:def-effect-generator}
    \G_{(q_0)}(s) := \frac{1}{sT} \sum_{q=1}^{q_0} \Phi_q(sT,c(sT)^{p+1}),
\end{align}
where the truncation order $q_0$ is determined by \thm{general-bch-truncation} so that error induced by the truncation, $\|\exp({\G_{(q_0)}(s)T})-\tilde{\mc P}(sT,c(sT)^{p+1})^{1/s}\|$, is within the required error tolerance.
The Richardson extrapolation error is therefore determined by the power series of the truncated scaled effective generator $\G_{(q_0)}(s)$.

To analyze the power series of $\G_{(q_0)}(s)$, we first show that each $\Phi_q(t,\lambda)$ is a homogeneous polynomial of $t$ and $\lambda$ and bound the coefficients of the polynomial. For ease of presentation, we also define the generators without coefficients as
\begin{align}
\label{eq:def-K-tilde}
 	\tilde{\mathcal{K}}_j := \begin{cases}\ad_{H_{\gamma_v}} & \text{if } j = 2v-1\\ \mathcal{L}_v & \text{if } j = 2v \end{cases}
\end{align}
for $v \in [M]$.
\begin{lemma}[BCH Coefficient Bound]
\label{lem:bch-bound}
 	Define the sum of the commutator norms for the base generators $\tilde{\mc K}_j$ as
 	\begin{align}
 	\label{eq:alpha-q-q}
 	\alpha_{\comm}^{(q_1, q_2)}=\sum_{(j_1, \ldots, j_q)\in S_{q_1, q_2}}\big\|[\tilde{\mc K}_{j_1}, \tilde{\mc K}_{j_2},\ldots,\tilde{\mc K}_{j_q}]\big\|,
 	\end{align}
 	where $q = q_1+q_2$ and
 	\begin{align}
 	\label{eq:S-q-q}
 	 	S_{q_1, q_2}:=\Bigl\{(j_1, \ldots, j_q)\in [2M]^q\Bigm| &|\{k\in [q]\mid j_k \text{ is odd}\}|=q_1, |\{k\in [q]\mid j_k \text{ is even}\}|=q_2\Bigr\}. \nonumber
 	\end{align}
 	The $q$-th order term $\Phi_q(t,\lambda)$ in the BCH formula of $\tilde{\P}(t,\lambda)$ can be written as a homogeneous polynomial of $t$ and $\lambda$:
 	\begin{align}
 	\Phi_q(t,\lambda) = \sum_{\substack{q_1, q_2 \ge 0\\q_1+q_2 = q}}\Phi_{q_1, q_2} t^{q_1}\lambda^{q_2},
 	\end{align}
 	such that the coefficients $\Phi_{q_1, q_2}$ are bounded by
 	\begin{align}
 	\|\Phi_{q_1, q_2}\|\le \frac{a_{\max}^{q_1}}{q^2}\alpha_{\comm}^{(q_1, q_2)}.
 	\end{align}
\end{lemma}
\begin{proof}
 	Substituting Eq.~\eq{def-K} into the expression of the $q$-th order BCH term $\Phi_q(t,\lambda)$ (Eq.~\eq{bch}) yields 
 	\begin{align}
 	 	\Phi_q(t,\lambda)
 	 	=~ \begin{multlined}[t]
 	 	 	\sum_{\substack{p_v,p'_v\ge 0\\ p_1+p_1'+\cdots+p_M+p_M'=q }}\bigg(\prod_{v=1}^M \frac{\lambda^{p_v}(-\i ta_v)^{p'_v}}{p_v!p'_v!} \bigg) \\ 
 	 	\cdot \sum_{\sigma\in S_q} c_{\sigma, q}R_{\sigma}\big\{[\underbrace{\mc L_M, \ldots,\mc L_M}_{p_M}, \underbrace{\A_{\gamma_M}, \ldots, \mc A_{\gamma_M}}_{p_M'}, \ldots, \underbrace{ \mc L_1,  \ldots, \mc L_1}_{p_1}, \underbrace{\A_{\gamma_1},  \ldots, \mc A_{\gamma_1}}_{p_1'}]\big\}.\label{eq:Z_q}\end{multlined}
 	\end{align}
 	
 	Note that any summand in Eq.~\eq{Z_q} containing $[\tilde{\mc K}_{j_1}, \tilde{\mc K}_{j_2}, \ldots, \tilde{\mc K}_{j_q}]$ for $(j_1, \ldots, j_q)\in [2M]^q$ is a monomial of $\lambda$ and $t$. The degree of $t$ is $|\{k\in [q]\mid j_k \text{ is odd}\}|=: q_1$ and the degree of $\lambda$ is $|\{k\in [q]\mid j_k \text{ is even}\}|=: q_2$. We now bound the coefficient of this monomial. The commutator $[\tilde{\mc K}_{j_1}, \tilde{\mc K}_{j_2}, \ldots, \tilde{\mc K}_{j_q}]$ can only originate from one summand in the sum over $p_v, p'_v$ in Eq.~\eq{Z_q} because the operator counts $(p_1, p_1', \ldots, p_M, p_M')$ are uniquely determined by $(j_1, \ldots, j_q)$. The number of permutations $\sigma\in S_q$ in the second sum in Eq.~\eq{Z_q} that preserve the commutator is
 	\begin{align*}
 	 	\prod_{v=1}^M p_v!p'_v!.
 	\end{align*}
 	Thus, the coefficient of the monomial $t^{q_1}\lambda^{q_2}[\tilde{\mc K}_{j_1}, \ldots, \tilde{\mc K}_{j_q}]$ can be upper bounded by
 	\begin{align*}
 	 	\bigg\|\frac{1}{q^2}\prod_{v=1}^M \frac{(-\i a_v)^{p'_v}}{p_v!p'_v!}\prod_{v=1}^M p_v!p'_v!\big\|[\tilde{\mc K}_{j_1}, \tilde{\mc K}_{j_2}, \ldots, \tilde{\mc K}_{j_q}]\big\|\bigg\| \le \frac{a_{\max}^{q_1}}{q^2}  \big\|[\tilde{\mc K}_{j_1}, \tilde{\mc K}_{j_2}, \ldots, \tilde{\mc K}_{j_q}]\big\|,
 	\end{align*}
 	as $|c_{\sigma, q}|\le 1/q^2$. Then, we regroup the sum in Eq.~\eq{Z_q} according to $(j_1, \ldots, j_q)\in S_{q_1, q_2}$ (which have the same degree of $t$ and $\lambda$). The resulting polynomial
 	\begin{align*}
 	 	\Phi_q(t,\lambda) = \sum_{\substack{q_1, q_2 \ge 0\\q_1+q_2 = q}}\Phi_{q_1, q_2} t^{q_1}\lambda^{q_2}
 	\end{align*}
 	satisfies
 	\begin{align*}
 	 	\|\Phi_{q_1, q_2}\|\le  \frac{a_{\max}^{q_1}}{q^2}\sum_{(j_1, \ldots, j_q)\in S_{q_1, q_2}} \big\|[\tilde{\mc K}_{j_1}, \tilde{\mc K}_{j_2}, \ldots, \tilde{\mc K}_{j_q}]\big\| = \frac{a_{\max}^{q_1}}{q^2}\alpha_{\comm}^{(q_1, q_2)}.
 	\end{align*}
\end{proof}


By \lem{bch-bound}, $\G_{(q_0)}(s)$ is a power series in $s$ and we denote the $q$-th order coeffcient by $\E_{(q_0),q+1}$ such that \begin{align}
\label{eq:def-effect-generator}
    \G_{(q_0)}(s)=\sum_{q=0}^{\infty} \E_{(q_0),q+1} (sT)^q. 
\end{align}
We bound the norm of $\E_{(q_0),q}$ in the following lemma. 
\begin{lemma}
    \label{lem:bound-E}
    Define a constant
    \begin{align}\label{eq:def-mu}
        \mu_{\mathrm{comm},q_0} = \max_{q' \ge p+1} \Bigg(\sum_{\substack{1\le q_1+ q_2\le q_0 \\
        q_1+(p+1)q_2 = q'}}\frac{a_{\max}^{q_1}c^{q_2}}{(q_1+q_2)^2} \alpha_{\rm comm}^{(q_1, q_2)}\Bigg)^{1/q'}.
    \end{align}
    Then we have \begin{align}
        \G_{(q_0)}(s) = -\i \ad_H + \sum_{q=p}^{\infty} \E_{(q_0),q+1} (sT)^q \label{eq:def-E},
    \end{align} where $\E_{(q_0),q}$ can be bounded by \begin{align}
        \|\E_{(q_0),q}\| \le  \mu_{\mathrm{comm},q_0}^q.
    \end{align}
\end{lemma}
\begin{proof}
    By the order condition of the product formula (see \citep[Lemma 2]{watson2025exponentially}), we have \begin{align*}
        \frac{1}{sT}\sum_{q=1}^{\infty} \Phi_{q}(sT,0) = -\i \ad_H + O((sT)^p),
    \end{align*}
    which implies \begin{align*}
        \Phi_{q_1, 0} = 0,
    \end{align*}
    for all $1 <q_1 \le p$, and \begin{align*}
        \Phi_{1,0} = - \i \ad_H sT.
    \end{align*} 
    Therefore, for any $c>0$, we have \begin{align*}
        \frac{1}{sT} \Phi_{q}(sT,c(sT)^{p+1})&= \frac{1}{sT}\Big(\sum_{q_1=1}^{q_0}\Phi_{q_1, 0}(sT)^{q_1} + \sum_{q = 1}^{q_0}\sum_{\substack{q_1\ge 0,q_2\ge 1\\ q_1+q_2 = q}} \Phi_{q_1, q_2} (sT)^{q_1}(c(sT)^{p+1})^{q_2}\Big) \\
        & = -\i \ad_H + \sum_{q_1 = p+1}^{q_0} \Phi_{q_1, 0} (sT)^{q_1-1} + O((sT)^{p}) \\
        & = -\i \ad_H + \sum_{q = p}^{\infty} \E_{(q_0),q+1} (sT)^{q}.
    \end{align*}
    For any $q\ge p+1$, we have \begin{align*}
        \|\E_{(q_0),q}\| &= \bigg\|\sum_{q'=1}^{q_0}\sum_{\substack{q_1, q_2\ge 0,~ q_1+q_2=q' \\
        q_1+(p+1)q_2 = q}} \Omega_{q_1,q_2} c^{q_2} \bigg\|\\
        &\le \sum_{\substack{1\le q_1+ q_2\le q_0 \\
        q_1+(p+1)q_2 = q}}\frac{a_{\max}^{q_1}c^{q_2}}{(q_1+q_2)^2} \alpha_{\rm comm}^{(q_1, q_2)} \\
        &\le  \mu_{\mathrm{comm},q_0}^q.
    \end{align*}
    %For $q=p+1$, we have \begin{align*}
    %    \|\E_{p+1}\| &\le c\Big\|\sum_{v=1}^M \L_v\Big\|+\|\Omega_{p+1,0}\| \le c\alpha_{\rm comm}^{(0,1)}+\frac{a_{\max}^{p+1}}{(p+1)^2}\alpha_{\rm comm}^{(p+1,0)}\le \mu_{\mathrm{comm},q_0}^{p+1.}
    %\end{align*}
\end{proof}
\begin{lemma}
\label{lem:expansion-exp-G(s)}
    The evolution $e^{\G_{(q_0)}(s) T}$ admits series expansion in $s$ given by \begin{align*}
        e^{\G_{(q_0)}(s) T}=e^{-\i \ad_H T}+ \sum_{q=p}^{\infty}\tilde{\E}_{(q_0),q} s^q,
    \end{align*}
    where $\tilde{\E}_{(q_0),q}$ are superoperators bounded by \begin{align*}
       \|\tilde{\E}_{(q_0),q}\|\le   e(2\mu_{\mathrm{comm},q_0}T)^{q(1+1/p)},
    \end{align*}
    if $2\mu_{\mathrm{comm},q_0}T\ge 1$, and bounded by \begin{align*}
        \|\tilde{\E}_{(q_0),q}\|\le e(2\mu_{\mathrm{comm},q_0}T)^{q},
    \end{align*}
    if $2\mu_{\mathrm{comm},q_0}T\le 1$.
\end{lemma}
\begin{proof}
By Dyson series expansion in the interaction picture, we have 
\begin{align*}
    &e^{\G_{(q_0)}(s) T}-e^{-\i \ad_H T} \\
    = \ & e^{-\i \ad_H T}  \sum_{j=1}^{\infty} \int_{0}^T \d \tau_1 \int_{0}^{\tau_1}\d \tau_2 \cdots \int_{0}^{\tau_{j-1}}\d \tau_j \prod_{j'=j}^{1}\bigg(e^{\i\ad_H \tau_{j'}}\bigg(\sum_{q=p}^{\infty}\E_{(q_0),q+1}s^qT^q\bigg)e^{-\i \ad_H \tau_{j'}}\bigg) \\
    = \ &  e^{-\i \ad_H T}\sum_{q=p}^{\infty}s^{q} \sum_{j=1}^{\infty} \sum_{\substack{q_1,\ldots,q_j \ge p \\ q_1+\cdots+q_j=q}} \int_{0}^T \d \tau_1 \int_{0}^{\tau_1}\d \tau_2 \cdots \int_{0}^{\tau_{j-1}}\d \tau_j \prod_{j'=j}^{1}\bigg(e^{\i\ad_H \tau_{j'}}\E_{(q_0),q_{j'}+1}e^{-\i \ad_H \tau_{j'}}T^{q_{j'}}\bigg)\\
    =:\,& \sum_{q=p}^{\infty}\tilde{\E}_{(q_0),q} s^q   
\end{align*}
The superoperator $\tilde{\E}_{(q_0),q}$ can be bounded by \begin{align*}
    \big\|\tilde{\E}_{(q_0),q}\big\|&\le \sum_{j=1}^{\infty} \sum_{\substack{q_1,\ldots,q_j \ge p \\ q_1+\cdots+q_j=q}} \int_{0}^T \d \tau_1 \int_{0}^{\tau_1}\d \tau_2 \cdots \int_{0}^{\tau_{j-1}}\d \tau_j \prod_{j'=1}^{j}\big(\big\|\E_{(q_0),q_{j'}+1}\big\|T^{q_{j'}}\big) \\
    &\le  \sum_{j=1}^{\infty} \sum_{\substack{q_1,\ldots,q_j \ge p \\ q_1+\cdots+q_j=q}}\frac{T^{q+j}}{j!}\prod_{j'=1}^{j}\mu_{\mathrm{comm},q_0}^{q_{j'}+1} \\
    &=\sum_{j=1}^{\lfloor q/p\rfloor }\frac{(\mu_{\mathrm{comm},q_0}T)^{q+j}}{j!}\Big(\sum_{\substack{q_1,\ldots,q_j \ge q \\ q_1+\cdots+q_j=q}} 1\Big).
\end{align*}
The second sum can be bounded by \begin{align*}
    \sum_{\substack{q_1,\ldots,q_j \ge q \\ q_1+\cdots+q_j=q}} 1 &\le \sum_{\substack{ q_1,\ldots,q_j \ge 0 \\ q_1+\cdots+q_j=q}} 1=\binom{q+j-1}{j-1} \le 2^{q+j}.
\end{align*}
Therefore, for $2\mu_{\mathrm{comm},q_0}T\ge 1$, we have \begin{align*}
    \big\|\tilde{\E}_{(q_0),q}\big\|\le \sum_{j=1}^{\lfloor q/p\rfloor }\frac{(2\mu_{\mathrm{comm},q_0}T)^{q+j}}{j!} 
    \le (2\mu_{\mathrm{comm},q_0}T)^{q+q/p}\sum_{j=1}^{\infty}\frac{1}{j!}\le e(2\mu_{\mathrm{comm},q_0}T)^{q(1+1/p)} ,
\end{align*}
and for $2\mu_{\mathrm{comm},q_0}T\le 1$, we have \begin{align*}
    \big\|\tilde{\E}_{(q_0),q}\big\| \le (2\mu_{\mathrm{comm},q_0}T)^{q}\sum_{j=0}^{\infty}\frac{ (2\mu_{\mathrm{comm},q_0}T)^{j}}{j!}\le (2\mu_{\mathrm{comm},q_0}T)^{q} e^{2\mu_{\mathrm{comm},q_0}T}\le e(2\mu_{\mathrm{comm},q_0}T)^{q}.
\end{align*}
\end{proof}
\xz{Also assume that the noise channel after each layer can be implemented efficiently. }
\begin{theorem}
\label{thm:general-application}
    Let $H = \sum_{\gamma = 1}^{\Gamma} H_{\gamma}$ be the target Hamiltonian. Suppose we approximate its evolution by implementing a $p$-th order product formula with $M$ gate layers:
    \begin{align*}
        P(t) = \prod_{v=1}^M e^{-\i t a_v H_{\gamma_v}},
    \end{align*}
    and that the implementation is subject to physical noise described by $e^{\lambda \L_v}$ after the $v$-th gate layer, with a tunable noise strength $\lambda \ge \lambda_{\min}$.
    
    Let $c>0$ and $q_0\in \mb Z_+$. Let $\bar{\beta}_{\mathrm{comm},q_0}(s)$ be the doubly nested commutator bound (Eq.~\eq{def-beta-comm-sum}) evaluated for the generators:
    \begin{align*}
        \mathcal{K}_j= \begin{cases}-\i sT a_{v}\ad_{H_{\gamma_v}} & \text{if } j = 2v-1\\ c(sT)^{p+1}\mathcal{L}_v & \text{if } j = 2v \end{cases}.
    \end{align*}
    Let $\mu_{\mathrm{comm},q_0}$ be the parameter defined in Eq.~\eq{def-mu}. 
    
    There exists an algorithm that, given any observable $O$ and initial state $\rho_0$, estimates $\tr[O e^{-\i \ad_H T}\rho_0]$ to within an additive error of $\epsilon\|O\|$ using a total of \begin{align*}
        O\Big( \max\{\log^3(1/\eps), \log^2(1/\eps)s_0^{-1}\}\frac{(\log\log (1/\eps))^3}{\eps^2}\Big)  = \widetilde{O}\Big(\frac{1}{s_0 \eps^2}\Big)
    \end{align*}Trotter steps. This involves $\widetilde{O}(1/\epsilon^2)$ independent circuit runs, each comprising at most 
    \begin{align*}
        O( \max\{\log^3(1/\eps), \log^2(1/\eps)s_0^{-1}\}) = \widetilde{O}\left( \frac{1}{s_0} \right)
    \end{align*} Trotter steps,
    where $s_0$ is the largest step size satisfying
    \begin{align} \label{eq:s0-general}
       s_0 (2\mu_{\mathrm{comm},q_0}T)^{1+1/p} \le 1/e \quad \text{and} \quad
       2e \bar{\beta}_{\mathrm{comm},q_0}(s_0) \le s_0 \frac{\epsilon}{4eC\log m},
    \end{align}
    with $C$ being the Richardson constant and $m = O(\log(1/\eps))$ satisfying $2e^{1-m} \le \eps /(4C\log m)$.
    
    This result holds provided that the native noise strength of the device, $\lambda_{\min}$, satisfies the feasibility condition $\lambda_{\min} = \widetilde{O}(c (s_0 T)^{p+1})$.
\end{theorem}
\begin{proof}
    Recall that the function for extrapolation is \begin{align*}
        f(s)&=\tr[O(\tilde{\mc P}(sT, c(sT)^{p+1}))^{1/s}\rho_0],
    \end{align*}
    for $s>0$ and $1/s\in \mb Z$.
    Here we consider the scenario where $2\mu_{\mathrm{comm},q_0} T \ge 1$. By \lem{expansion-exp-G(s)}, the evolution $e^{\G_{(q_0)}(s) T}$ admits series expansion in $s$ given by \begin{align*}
        e^{\G_{(q_0)}(s) T}=e^{-\i \ad_H T}+ \sum_{q=p}^{\infty}\tilde{\E}_{(q_0),q} s^q,
    \end{align*}
    such that \begin{align*}
        \|\E_{(q_0), q}\|\le e (2\mu_{\mathrm{comm},q_0}T)^{q(1+1/p)}.
    \end{align*}    
    Then the function $f(s)$ can be decomposed as \begin{align*}
        f(s)&= \tr[Oe^{\G_{(q_0)}(s)T} \rho_0]+\tr[O\bigl((\tilde{\mc P}(sT, c(sT)^{p+1}))^{1/s}-e^{\G_{(q_0)}(s)T}\bigr)\rho_0] \\
        &= \underbrace{\tr[Oe^{-\i \ad_H T}\rho_0]+\sum_{q=p}^{m-1}\tr[O \E_{(q_0), q}\rho_0] s^q}_{P_m(s)}+\underbrace{\sum_{q=m}^{\infty}\tr[O\E_{(q_0), q}\rho_0] + \tr[O(\tilde{\mc P}(sT, c(sT)^{p+1})-e^{\G_{(q_0)}(s)T})\rho_0]}_{R_m(s)},
    \end{align*}
    where $P_m$ is a degree-$(m-1)$ polynomial in $s$. Let $s_1, \ldots, s_m$ and $b_1, \ldots, b_m$ be the extrapolation points and coefficients defined in \lem{richardson-extrapolation} with $t=1$ and $s = s_0$ to be determined later, and then we have $1\le\|\mathbf{b}\|\le C\log m$, $1/s_j\in \mathbb{Z}$, and $1/s_j = r_j = O(\max\{m^3, m^2/s\}/j^2)$. We can set \begin{align*}
        m = O(\log(1/\eps)),
    \end{align*}
    such that \begin{align}
    \label{eq:cond-m}
        2e^{1-m} \le \frac{\eps}{4C\log m}.
    \end{align}
    
   
    For any $s$ satisfying \begin{align}
    \label{eq:cond-s}
        e\bar\beta_{\mathrm{comm},q_0}(s)\le s\frac{\eps}{4eC \log m} \le s \quad s(2\mu_{\mathrm{comm},q_0}T)^{1+1/p} \le 1/e,
    \end{align} the remainder $R_m(s)$ is bounded as \begin{align*}
        &|R_m(s)| \\
        \le ~&\sum_{q=m}^{\infty}\|O\|\|\tilde{\E}_{(q_0),q}\|\|\rho_0\|_1 s^q + \|O\|\Big\|(\tilde{\mc P}(sT, c(sT)^{p+1}))^{1/s}-(e^{\G_{(q_0)}(s)sT})^{1/s}\Big\|\|\rho_0\|_1 \\
        \le ~& \sum_{q=m}^{\infty} e(s(2\mu_{\mathrm{comm},q_0} T)^{1+1/p})^{q}\|O\|+\frac{(\max\{1, \|e^{\G_{(q_0)}(s) sT}\|\})^{1/s}}{s}\Bigl\|\tilde{\mc P}(sT, c(sT)^{p+1})-e^{\G_{(q_0)}(s)sT}\Bigr\|\|O\|\\
        \le ~&2e^{1-m} \|O\|+ \frac{(1+e\beta_{\mathrm{comm},q_0}(s))^{1/s}}{s}e\beta_{\mathrm{comm},q_0}(s)\|O\| \\
        \le ~& \frac{\eps}{4C\log m} \|O\|+(1+s)^{1/s}\frac{\eps}{4eC\log m}\|O\| \le \frac{\eps}{2C\log m} \|O\|
    \end{align*}
    where the third line follows from \lem{expansion-exp-G(s)} and \lem{power-bound}, the fourth line follows from 
    \thm{general-bch-truncation}, and the fifth line follows from Eq.~\eq{cond-m} and Eq.~\eq{cond-s}.
    
    Now, we bound the error of the extrapolation function \begin{align*}
        F^{(m)} = \sum_{j=1}^m b_j f(s_j).
    \end{align*}
    Set \begin{align*}
        f(0) :=\lim_{s\to 0} f(s) = \tr[Oe^{-\i \ad_H T}\rho_0].
    \end{align*} Let $s_0$ be the largest number smaller than $e^{-1}(2\mu_{\mathrm{comm},q_0}T)^{-1-1/p}$ and satisfying \begin{align*}
        e\bar\beta_{\mathrm{comm},q_0}(s_0)\le s_0\frac{\eps}{4eC \log m}
    \end{align*}
    Observe that $\bar\beta_{\mathrm{comm}, q_0}(s)$ consists of sums of norms of nested commutators involving $-\i sT a_v \mc{A}_{\gamma_v}$ and $c(sT)^{p+1} \mc L_{v}$ with grade $q \ge q_0+1$. Consequently, for $s \ge 0$, it is a polynomial in $s$ with non-negative coefficients and a lowest degree of $q_0+1$. This implies that the function $\bar\beta_{\mathrm{comm}, q_0}(s)/s$ is monotonically increasing for $s \ge 0$. 
    By \lem{richardson-extrapolation}, the extrapolation nodes satisfy $s_j \le s_0$. Therefore, as the condition in Eq.~\eq{cond-s} holds for the base step size $s_0$, it is satisfied for all $s_j$.
    Then we have \begin{align*}
        |F^{(m)}(s)-f(0)|&\le \|\mathbf{b}\|_1 \max_j|R_m(s_j)| \\
        &\le C\log m\frac{\eps}{2C\log m}\|O\|\\
        &\le \frac{\eps }{2}\|O\|.
    \end{align*}
    To estimate $f(s_j)$ with error $\eps\|O\|/(2\|\mathbf{b}\|_1)$ with success probability $1-1/(3m)$, we need to repeatedly run $(\tilde{\mc P}(s_jT, c(s_jT)^{p+1}))^{1/s_j}$ $O(\|\mathbf{b}\|_1^2/\eps^2\log m)$ times.  Each application of $(\tilde{\mc P}(s_jT, c(s_jT)^{p+1}))^{1/s_j}$ requires $1/s_j$ Trotter steps. Hence the maximum number of Trotter steps per coherent run is \begin{align*}
        O(\max_j r_j) & = O( \max\{m^3, m^2s_0^{-1}\})\\
        &=O( \max\{\log^3(1/\eps), \log^2(1/\eps)s_0^{-1}\}) \\
        &= \widetilde{O}\Big(\frac{1}{s_0}\Big),
    \end{align*} and the total number of Trotter steps is \begin{align*}
        O\Big(\sum_{j=1}^m r_j\frac{\|\mathbf{b}\|_1^2}{\eps^2}\log m\Big) & = O\Big(\sum_j^{m} \max\{m^3, m^2s_0^{-1}\}\frac{1}{j^2}\frac{\log^3 m}{\eps^2}\Big)\\
        &=O\Big( \max\{\log^3(1/\eps), \log^2(1/\eps)s_0^{-1}\}\frac{(\log\log (1/\eps))^3}{\eps^2}\Big) \\
        &= \widetilde{O}\Big(\frac{1}{s_0\eps^2}\Big)
    \end{align*}
    The smallest noise strength is \begin{align*}
        \min_{j}c(s_jT)^{p+1} &= \Omega\Big(\min_{j}c\Big(\frac{j^2T}{\max\{m^3, m^2s_0^{-1}\}}\Big)^{p+1}\Big) \\
        &=\Omega\Big(c\min\Big\{\frac{T^{p+1}}{\log^{3p+3}(1/\eps)}, \frac{ T^{p+1}s_0^{p+1}}{\log^{2p+2}(1/\eps)}\Big\}\Big)\\
        & = \widetilde{\Omega}(c (s_0 T)^{p+1}).
    \end{align*}
\end{proof}
\xz{TBD: explaining three tunable parameters: $c$ balances Trotter error and physical error, $q_0$ balances BCH truncation error and extrapolation error.}

\section{Applications to $k$-local Hamiltonians}
\label{sec:local}
Suppose that $\ad_H=\sum_{\gamma\in[\Gamma]}\A_\gamma$ and $\mc L_{\tot}:=\sum_{v\in [M]}\mc L_v$ are $k$-local Lindbladians on a lattice $\Lambda = [N]$. We further assume that $\ad_H$ is $g_H$-extensiven and $\L_{\rm tot}$ is $g_L$-extensive.
Recall the unified representation of $\mc A_{\gamma}$ and $\mc L_v$ in Eq.~\eq{def-K} and Eq.~\eq{def-K-tilde}: \begin{align*}
    \mathcal{K}_j = \begin{cases}-\i ta_{v}\ad_{H_{\gamma_v}} & \text{if } j = 2v-1 \\ \lambda\mathcal{L}_v & \text{if } j = 2v \end{cases}, \quad \tilde{\mathcal{K}}_j := \begin{cases}\ad_{H_{\gamma_v}} & \text{if } j = 2v-1\\ \mathcal{L}_v & \text{if } j = 2v \end{cases} 
\end{align*}
for $v=1, 2,\ldots, M$.  

Now we substitute $t = sT, \lambda = c (sT)^{p+1}$ and let $\G_{(q_0)}(s)$ be the effective generator defined in Eq.~\eq{def-effect-generator}. By \lem{bound-E}, the coefficients of the series expansion of $e^{\G_{(q_0)}(s)T}$ can be bounded as powers of $\mu_{\mathrm{comm},q_0}$ defined in Eq.~\eq{def-mu}. We now provide an upper bound of $\mu_{\mathrm{comm},q_0}$. 
\begin{lemma}
\label{lem:bound-mu}
    The parameter $\mu_{\mathrm{comm},q_0}$ defined in Eq.~\eq{def-mu} satisfies \begin{align*}
        \mu_{\mathrm{comm},q_0} \le q_0 N^{1/(p+1)}(2a_{\max}k \Upsilon g_H+ (2ckg_L)^{1/(p+1)}).
    \end{align*}
\end{lemma}
\begin{proof}
    For the set $S_{q_1, q_2}$ defined in Eq.~\eq{alpha-q-q}, we decompose it according to the position of the odd and even indices.  For any index set $\mathcal{I} \subseteq [q_1+q_2]$, we define the subset $S_{\mathcal{I}}$ as
    \[
    S_{\mathcal{I}} = \Bigl\{ (j_1, \ldots, j_{q_1+q_2}) \in [2M]^{q_1+q_2} \Bigm| 
        j_k \text{ is odd for } k \in \mathcal{I}, \text{ and } 
        j_k \text{ is even for } k \in [q_1+q_2] \setminus \mathcal{I}
    \Bigr\}.
    \]
    We can then express $S_{q_1, q_2}$ as the disjoint union over the collection of all size-$q_1$ subsets of $[q_1+q_2]$:
    \[
    S_{q_1, q_2} = \bigsqcup_{\mathcal{I} \in \binom{[q_1+q_2]}{q_1}} S_{\mathcal{I}}.
    \]
    The commutator bounds are \begin{align*}
    \alpha_{\rm comm}^{(q_1,q_2)} &=  \sum_{(j_1, \ldots, j_{q_1+q_2})\in S_{q_1, q_2}}\big\|[\tilde{\mc K}_{j_1}, \tilde{\mc K}_{j_2},\ldots,\tilde{\mc K}_{j_{q_1+q_2}}]\big\|\\
    &= \sum_{\mc I\in \binom{[q_1+q_2]}{q_1}}\sum_{(j_1, \ldots, j_{q_1+q_2})\in S_{\mc I}}\big\|[\tilde{\mc K}_{j_1}, \tilde{\mc K}_{j_2},\ldots,\tilde{\mc K}_{j_{q_1+q_2}}]\big\|\\
    &\le \sum_{\mc I\in \binom{[q_1+q_2]}{q_1}} N(q_1+q_2-1)!(2k)^{q_1+q_2-1}(\Upsilon g_H)^{q_1}g_L^{q_2}\\
    &= \binom{q_1+q_2}{q_1}N(q_1+q_2-1)!(2k)^{q_1+q_2-1}(\Upsilon g_H)^{q_1}g_L^{q_2}\\
    &\le \binom{q_1+(p+1)q_2}{q_1}N(q_1+q_2)!(2k\Upsilon g_H)^{q_1}((2kg_L)^{1/(p+1)})^{q_2(p+1)},
    \end{align*}
    where the third line follows from \cor{commu-bound-sum}. Then we have  \begin{align*}
        \mu_{\mathrm{comm},q_0} &= \max_{q' \ge p+1} \Bigg(\sum_{\substack{1\le q_1+q_2\le q_0 \\
        q_1+(p+1)q_2 = q'}}\frac{a_{\max}^{q_1}c^{q_2}}{(q_1+q_2)^2} \alpha_{\rm comm}^{(q_1, q_2)}\Bigg)^{1/q'} \\
        &\le \max_{q' \ge p+1} \bigg(\sum_{\substack{1\le q_1+q_2\le q_0 \\
        q_1+q_2(p+1) = q'}}\binom{q_1+q_2(p+1)}{q_1}N(q_1+q_2)!(2a_{\max}k\Upsilon g_H)^{q_1}((2ckg_L)^{1/(p+1)})^{q_2(p+1)}\bigg)^{1/q'}\\
        &\le \max_{q' \ge p+1} \bigg(\sum_{q_1=1}^{q_0}\binom{q'}{q_1}Nq_0^{q'}(2a_{\max}k\Upsilon g_H)^{q_1}((2ckg_L)^{1/(p+1)})^{q'-q_1}\bigg)^{1/q'}\\
        &\le q_0N^{1/(p+1)}\max_{q'\ge p+1} \big((2a_{\max}k\Upsilon g_H+(2ckg_L)^{1/(p+1)})^{q'}\big)^{1/q'}\\
        &=q_0N^{1/(p+1)}(2a_{\max}k\Upsilon g_H+(2ckg_L)^{1/(p+1)}).
    \end{align*}
\end{proof}

Now we apply \thm{general-application} to $k$-local Hamiltonians.  
\begin{theorem}
\label{thm:extrapolation-k-local}
    Suppose that the native noise strength of the device, $\lambda_{\min}$, satisfies
    \begin{align}
    \label{eq:cond-lambda}
        \lambda_{\min} = \widetilde{O}\left(\frac{1}{(Nk\Upsilon a_{\max}g_HT)^{1+1/p} kg_L}\right).
    \end{align} 
    There exists an algorithm that, given any observable $O$ and initial state $\rho_0$, estimates $\tr[O e^{-\i \ad_H T}\rho_0]$ to within an additive error of $\epsilon\|O\|$ using a total of 
    \begin{align*}
        O\left( \max\left\{\log^3(1/\eps), \Xi \log^2(1/\eps)\right\} \frac{(\log\log (1/\eps))^3}{\eps^2} \right)
    \end{align*}
    Trotter steps, where we define the parameter dependent factor
    \begin{align*}
        \Xi := N^{1/p}(k\Upsilon a_{\max} g_H T)^{1+1/p} \log^{1+1/p}(N/\eps).
    \end{align*}
    This involves $\widetilde{O}(1/\epsilon^2)$ independent circuit runs, each comprising at most
    \begin{align*}
        O\left( \max\left\{\log^3(1/\eps), \Xi \log^2(1/\eps)\right\} \right)
    \end{align*}
    Trotter steps.
\end{theorem}

\begin{proof}
    Note that $\mc K:= \sum_{j=1}^{2M} {\mc K}_j$ is $k$-local and $(\Upsilon sT a_{\max}g_H+c(sT)^{p+1} g_L)$-extensive.
    Following the proof of \lem{bch-approx}, the doubly right-nested commutator bound of $\mc K$, denoted $\bar{\beta}_{\mathrm{comm}, q_0}(s)$, satisfies 
    \begin{align*}
        \bar{\beta}_{\mathrm{comm}, q_0}(s) \le 2(4eq_0 k(\Upsilon sT a_{\max}g_H+c(sT)^{p+1} g_L))^{q_0+1}N,
    \end{align*}
    provided $4eq_0 k(\Upsilon sT a_{\max}g_H+c(sT)^{p+1} g_L) \le 1/2$. 
    Let $m=O(\log(1/\eps))$ be the number of extrapolation points. We choose the expansion order as $q_0 = \lceil \log(8e^2N\log m/\eps) \rceil = O(\log(N/\eps))$. 
    
    We set the step size $s_0$ as
    \begin{align}
    \label{eq:s0-def}
        s_0 &= \frac{1}{(2eq_0 N^{1/(p+1)}T (4ek \Upsilon a_{\max} g_H + (4ekcg_L )^{1/(p+1)}))^{1+1/p}} \\
        & = \Theta\left( \frac{1}{ \log^{1+1/p}(N/\eps)  N^{1/p}T^{1+1/p}(k\Upsilon a_{\max} g_H +(kcg_L)^{1/(p+1)})^{1+1/p} } \right).\nonumber
    \end{align}
    One can verify that this choice of $s_0$ satisfies the condition $4eq_0 k(\Upsilon s_0T a_{\max}g_H+c(s_0T)^{p+1} g_L) \le 1/2$.
    Furthermore, we have
    \begin{align*}
         \frac{\bar{\beta}_{\mathrm{comm}, q_0}(s_0)}{s_0} & \le 2\max\{(s_0^{p/(p+1)}4eq_0k\Upsilon T a_{\max} g_H)^{q_0+1}, (4eq_0kcg_L s_0 ^pT^{p+1})^{q_0+1}\}N \\
         &\le 2e^{-q_0-1}N \le \frac{\eps}{8e^2C \log m}.
    \end{align*}
    Also, $s_0$ satisfies the condition related to $\mu_{\mathrm{comm}, q_0}$:
    \begin{align*}
        s_0 &\le \frac{1}{e(2q_0 N^{1/(p+1)} T(2 a_{\max} k\Upsilon g_H + (2ckg_L)^{1/(p+1)}))^{1+1/p}} \le \frac{1}{e(2\mu_{\mathrm{comm}, q_0}T)^{1+1/p}}. 
    \end{align*}
    
    To optimize the constraint on $\lambda_{\min}$ without increasing the asymptotic complexity, we set
    \begin{align}
    \label{eq:def-c}
        c = \frac{(k\Upsilon a_{\max}g_H)^{p+1}}{kg_L},
    \end{align}
    which implies $(kcg_L)^{1/(p+1)} = k\Upsilon a_{\max}g_H$. Substituting this into Eq.~\eqref{eq:s0-def}, the inverse step size scales as:
    \begin{align*}
        \frac{1}{s_0} = O\left( N^{1/p}(k\Upsilon a_{\max} g_H T)^{1+1/p} \log^{1+1/p}(N/\eps) \right) = O(\Xi).
    \end{align*}

    Now, applying \thm{general-application}, the algorithm requires a total number of Trotter steps of:
    \begin{align*}
         O\left( \max\left\{\log^3(1/\eps), \Xi \log^2(1/\eps)\right\} \frac{(\log\log (1/\eps))^3}{\eps^2} \right).
    \end{align*}
    The maximum number of Trotter steps per run is
    \begin{align*}
  O\left( \max\left\{\log^3(1/\eps), \Xi \log^2(1/\eps)\right\} \right).
    \end{align*}
    This result holds provided that the native noise strength satisfies:
    \begin{align*}
        \lambda_{\min} = \widetilde{O}(c(s_0T)^{p+1}) = \widetilde{O}\Big(\frac{1}{ (Nk\Upsilon a_{\max}g_HT)^{1+1/p} kg_L }\Big).
    \end{align*}
\end{proof}
\xz{This reaches the limit of error mitigation under Pauli noise \cite{wang2021noise, cai2021multi}.}
In the noiseless limit ($g_L = 0$), \thm{extrapolation-k-local} provides a resource analysis of the Trotter error mitigation algorithm proposed by \citet{watson2025exponentially} for $k$-local lattice Hamiltonians. This result parallels that of \citet{mizuta2025commutator}, which analyzed the complexity of multi-product formulas without physical noise using a different proof technique.

A crucial implication of \thm{extrapolation-k-local} is that, even when the device is subject to local noise satisfying Eq.~\eq{cond-lambda}, the noiseless extrapolation algorithm can be applied with only a constant complexity overhead by properly coupling the Trotter step size and the noise strength in the extrapolation. This implies that physical noise mitigation is achieved at virtually no asymptotic cost. Specifically, our choice of $c$ in Eq.~\eq{def-c} merely rescales the step size $s_0$ by a factor of $2^{-1-1/p}$. Consequently, the total complexity, which is proportional to $s_0^{-1}$, increases only by a constant factor of $2^{1+1/p}$.

\xz{I'm thinking whether to explicitly compare our joint mitigation algorithm with the nested mitigation one. Since the asymptotic complexity advantage of our algorithm is only $\log(1/\eps)(\log\log(1/\eps))^2$. If not, we can instead claim that our algorithm can mitigate the physical error almost for free while doing Trotter error mitigation. }


\section{Applications to geometrically local Hamiltonian}
\label{sec:geo-local}
In this section, we consider a $d$-dimensional lattice $\Lambda \subseteq \mathbb{Z}^d$ and assume that $\ad_H = \sum_{\gamma\in [\Gamma]} \A_{\gamma}$ is geometrical local. 
For any $x\in \Lambda$ and a positive number $R$, define $B_R(x)$ as the $\ell_1$ ball centered on $x$ with radius $R$: \begin{align*}
    B_R(x) = \{ y\in \Lambda:d(x,y) \le R\},
\end{align*}
where \begin{align*}
    d(x,y) = \sum_{i=1}^d  |x_i-y_i|
\end{align*}
is the $\ell_1$ distance between $x = (x_1, \ldots, x_d)$ and $y= (y_1, \ldots, y_d)$. 
Assume that we can group the summand in the decomposition of $\ad_H$ and $\L_{\rm tot}$ as \begin{align*}
    &\ad_H = \sum_{x\in {\Lambda}} \A_x, \quad \A_x = \sum_{\gamma\in \mc I_x} \A_{\gamma} 
\end{align*} where $\{\mc I_x\}_{x\in \Lambda}$ is a partition of $[\Gamma]$, such that \begin{align*}
    &\supp(\A_{\gamma}) \subseteq B_R(x), \quad \forall \gamma\in \mc I_x
\end{align*} for some positive number $R$. We also assume that the observable of interest $O$ is supported on $B_{R_O}(x_O)$ for some positive $R_O$ and $x_O\in \Lambda$. Define \begin{align*}
    J_H = \max_{x\in \Lambda}\sum_{\gamma\in \mc I_x} \|\A_{\gamma}\|.
\end{align*} Define $\Lambda_d = 2^d/d!$ such that \begin{align*}
    |B_R(x)| = \Lambda_dR^d + O(R^{d-1})\le C_d R^d, \quad \forall x\in \Lambda,
\end{align*}
for some $C_d\le 2\Lambda_d$.
Then $\ad_H$ and $\L_{\rm tot}$ are $g_H$ and $g_L$-extensive with \begin{align*}
    g_H &\le \max_{x\in \Lambda} \sum_{y\in B_R(x)}\sum_{\gamma\in \mc I_y} \|\A_{\gamma}\|\le |B_R(x)|J_H  = O(C_d R^d J_H).
\end{align*}
For any $l \ge 0$, define\begin{align*}
    \ad_{H_{(l)}} = \sum_{y\in B_{R_O+l+R}(x_O)} \A_y,
\end{align*}
be the truncated Hamiltonian terms within a ball around the support of $O$. The difference between the time-evolved observables $O$ under $\ad_{H_{(l)}}$ and $ \ad_H$ can be bounded as follows. 
\begin{lemma}[{\cite[Lemma 2]{rao2025stability}}]
    For any $l\ge 0$, we have \begin{align*}
        \left\| e^{\i \ad_H t}O - e^{\i \ad_{H_{(l)}}t}O\right\| \le |\supp(O)|\|O\|\min\{e^{-l/R}(e^{eC_dR^d t}-1), 1\}.
    \end{align*}
\end{lemma}
We can choose \begin{align*}
    l = eC_d R^{d+1}T+ R\log(2|\supp(O)|/\eps),
\end{align*}
such that for any initial state $\rho_0$ \begin{align*}
    \left| \tr[Oe^{-\i \ad_H T}\rho_0]- \tr[Oe^{-\i \ad_{H_{(l)}} T}\rho_0]\right| \le \left\| e^{\i \ad_HT}O - e^{\i \ad_{H_{(l)}}T}O\right\| \le \eps/2.
\end{align*}
Therefore, we can simulate the truncated Hamiltonian $H_{(l)}$ on $B_{R_0+l+2R}(x_O)$,  and the noise terms are restricted to $B_{R_O+l+2R}(x_O)$. Define
\begin{align*}
     \L_{(l), \rm tot} =  \sum_{v\in \in B_{R_O+l+2R}(x_O)} \L_v,
\end{align*}
which includes all possible  noise terms. 
\xz{TBD}
\bibliographystyle{plainnat}
\bibliography{ref}
\appendix

\section{Proof of \lem{doubly-nested-bound-general}}
\label{append:sec:proof-doubly-nested}

For clarity, we restate the lemma here.

\begin{lemma}[Restatement of \lem{doubly-nested-bound-general}] 
For any sequence of positive integers $q_1, \ldots, q_{\ell+1}$, the integrated norm of the corresponding term in the modified generator $\mc L_{(q_0)}$ is bounded as
 \begin{align*}
  \int_{0}^{M} \bigg\|\prod_{\ell'=\ell}^1 \ad_{\Omega_{q_{\ell'}}(\tau)}\cdot\dot\Omega_{q_{\ell+1}}(\tau)\bigg\| \,\d\tau \le \bigg(\prod_{i=1}^{\ell+1} \frac{1}{q_i^2}\bigg) \bar{\beta}_{\comm}^{(q_1, \ldots, q_{\ell+1})}.
 \end{align*}
\end{lemma}
\begin{proof}
Let $T_q(t) := \{(\tau_1, \ldots, \tau_q) \mid t \ge \tau_1 \ge \cdots \ge \tau_q \ge 0\}$ be the time-ordered $q$-simplex.
 The term to be bounded is
\begin{align*}
    \mc{X}_{\ell}(\tau) := [\Omega_{q_1}(\tau), \ldots, \Omega_{q_\ell}(\tau), \dot{\Omega}_{q_{\ell+1}}(\tau)].
\end{align*}
By substituting the integral definitions, using the multilinearity of the commutator, and applying the triangle inequality, we obtain:
\begin{align}
    \|\mc{X}_{\ell}(\tau)\| 
    &\le \begin{multlined}[t]\sum_{\substack{i = 1 , \ldots, \ell+1\\\sigma_i\in S_{q_{i}}}} \bigg(\prod_{i=1}^{\ell+1} |c_{\sigma_i, q_i}|\bigg) \int_{T_{q_1}(\tau)} \d\tau_{1,1}, \ldots,\d\tau_{1,q_1}\cdots\int_{T_{q_{\ell+1}-1}(\tau)}\d\tau_{\ell+1,1},\ldots,\d\tau_{\ell+1,q_{\ell+1}-1} \\
        \bigg\| \Big[ \mc R_{\sigma_1}\big\{[\mc L(\tau_{1,1}), \ldots, \mc L(\tau_{1,q_1})]\big\}, \ldots, \mc R_{\sigma_\ell}\big\{[\mc L(\tau_{\ell,1}), \ldots, \mc L(\tau_{\ell,q_\ell})]\big\}, \\
        \mc R_{\sigma_{\ell+1}}\big\{[\mc L(\tau), \mc L(\tau_{\ell+1,1}), \ldots, \mc L(\tau_{\ell+1,q_{\ell+1}-1})]\big\} \Big] \bigg\|\nonumber
    \end{multlined}\\
    & \le \begin{multlined}[t]\sum_{\substack{i = 1 , \ldots, \ell+1\\\sigma_i\in S_{q_{i}}}} \bigg(\prod_{i=1}^{\ell+1} \frac{1}{q_i^2} \frac{1}{q_i!}\bigg)q_{\ell+1} \int_{[\tau]^{q_1}} \d\tau_{1,1}, \ldots,\d\tau_{1,q_1}\cdots\int_{[\tau]^{q_{\ell+1}-1}}\d\tau_{\ell+1,1},\ldots,\d\tau_{\ell+1,q_{\ell+1}-1} \\
        \bigg\| \Big[ \mc R_{\sigma_1}\circ \mc T\big\{[\mc L(\tau_{1,1}), \ldots, \mc L(\tau_{1,q_1})]\big\}, \ldots, \mc R_{\sigma_\ell}\circ \mc T\big\{[\mc L(\tau_{\ell,1}), \ldots, \mc L(\tau_{\ell,q_\ell})]\big\}, \\
        \mc R_{\sigma_{\ell+1}}\circ \mc T\big\{[\mc L(\tau), \mc L(\tau_{\ell+1,1}), \ldots, \mc L(\tau_{\ell+1,q_{\ell+1}-1})]\big\} \Big] \bigg\|\nonumber
    \end{multlined}\\
    & \le \begin{multlined}[t]\sum_{\substack{i = 1 , \ldots, \ell+1\\\sigma_i\in S_{q_{i}}}} \bigg(\prod_{i=1}^{\ell+1} \frac{1}{q_i^2} \frac{1}{q_i!}\bigg)q_{\ell+1}
        \sum_{\substack{
             v_{1,1}, \ldots, v_{\ell, q_\ell} = 1 \\
             v_{\ell+1, 1}, \ldots, v_{\ell+1, q_{\ell+1}-1} = 1
            }}^{\lceil\tau\rceil} \bigg\| \Big[ \mc R_{\sigma_1}\circ \mc T\big\{[\mc L(v_{1,1}), \ldots, \mc L(v_{1,q_1})]\big\}, \ldots, \\
        \mc R_{\sigma_\ell}\circ \mc T\big\{[\mc L(v_{\ell,1}), \ldots, \mc L(v_{\ell,q_\ell})]\big\}, 
        \mc R_{\sigma_{\ell+1}}\circ \mc T\big\{[\mc L(\tau), \mc L(v_{\ell+1,1}), \ldots, \mc L(v_{\ell+1,q_{\ell+1}-1})]\big\} \Big] \bigg\|,\label{eq:bound-Xk-general}
    \end{multlined} 
\end{align}
where the third line follows from the same integral decomposition argument in the proof of \lem{magnus-bch-general}. For each $i\in[\ell]$, after applying time-ordering operator $\mc T$, each unordered commutator $[\mc L(v_{i,1}), \ldots, \mc L(v_{i,q_i})]$ is mapped to an ordered commutator 
\begin{align}
\label{eq:ordered-comm-1-general}
    [\underbrace{\mc L(\lceil\tau\rceil) \ldots,\mc L(\lceil\tau\rceil)}_{q_{i,\lceil\tau\rceil}},  \ldots, \underbrace{\mc L(1),  \ldots, \mc L(1)}_{q_{i,1}}] = [\underbrace{\mc K_{\lceil\tau\rceil} \ldots,\mc K_{\lceil\tau\rceil}}_{q_{i,\lceil\tau\rceil}},  \ldots, \underbrace{\mc K_{1},  \ldots, \mc K_{1}}_{q_{i,1}}],
\end{align}
where $q_{i,j}$ is the number of $j$ in $v_{i,1}, \ldots, v_{i,q_i}$ and there are
\begin{align*}
    \binom{q_i}{q_{i, \lceil\tau\rceil}, \ldots, q_{i,1}} = \frac{q_i!}{q_{i, \lceil\tau\rceil}! \cdots, q_{i,1}!}
\end{align*}
unordered commutators mapped to the same  ordered commutator in Eq.~\eq{ordered-comm-1-general}. Then, after applying $\mc R_{\sigma_i}$ to each ordered commutator, we regroup the sum according to the index sequence $(u_{i,1}, \ldots, u_{i,q_i})$ of the resulting unordered commutator $[\mc K_{u_{i,1}},\ldots, \mc K_{u_{i,q_i}}]$. Each unordered commutator $[\mc K_{u_{i,1}},\ldots, \mc K_{u_{i,q_i}}]$ is mapped from one ordered commutator $[\underbrace{\mc K_{\lceil\tau\rceil} \ldots,\mc K_{\lceil\tau\rceil}}_{q_{i,\lceil\tau\rceil}},  \ldots, \underbrace{\mc K_{1},  \ldots, \mc K_{1}}_{q_{i,1}}]$ and the number of $\sigma_i\in S_{q_i}$ mapping the ordered commutator to the same unordered commutator is
\begin{align*}
    q_{i, \lceil\tau\rceil}! \cdots, q_{i,1}!.
\end{align*}
Therefore, the combinatorial coefficient of $[\mc K_{u_{i,1}},\ldots, \mc K_{u_{i,q_i}}]$ in the final sum is 
\begin{align*}
    \frac{q_i!}{q_{i, \lceil\tau\rceil}! \cdots, q_{i,1}!}q_{i, \lceil\tau\rceil}! \cdots, q_{i,1}! = q_i!.
\end{align*}
For $i= \ell+1$, after applying $\mc T$, each $[\mc L(\tau), \mc L(v_{\ell+1,1}), \ldots, \mc L(v_{\ell+1,q_{\ell+1}-1})]$ is mapped to
\begin{align}
\label{eq:ordered-comm-general}
    [\mc K_{\lceil \tau \rceil},\underbrace{\mc K_{\lceil\tau\rceil} \ldots,\mc K_{\lceil\tau\rceil}}_{q_{\ell,\lceil\tau\rceil}},  \ldots, \underbrace{\mc K_{1},  \ldots, \mc K_{1}}_{q_{\ell,1}}],
\end{align}
where $q_{\ell,j}$ is the number of $j$ in $v_{\ell+1, 1},\ldots, v_{\ell+1, q_{\ell+1}-1}$, and there are
\begin{align*}
    \binom{q_{\ell+1}-1}{q_{i,\lceil \tau \rceil}, \ldots, q_{i,1}} = \frac{(q_{\ell+1}-1)!}{q_{\ell+1,\lceil \tau \rceil}!, \ldots, q_{\ell+1,1}!}
\end{align*} unordered commutators mapped to the same ordered commutator in Eq.~\eq{ordered-comm-general}. After applying $\mc R_{\sigma_{\ell+1}}$, we regrouping the sum of the commutators according to the position of the fixed $\mc K_{\lceil \tau \rceil}$ after permutation, $j'$, and the indices of other $\mc K$ operators in a commutator
\begin{align*}
    [\mc K_{v_{\ell+1,1}}, \ldots,\mc K_{\lceil\tau\rceil}, \ldots, \mc K_{v_{\ell+1,q_{\ell+1}-1}}].
\end{align*}
Each such unordered commutator is mapped from one ordered commutator in Eq.~\eq{ordered-comm-general} and the number of $\sigma_{\ell+1}\in S_{q_{\ell+1}}$ satisfying $\sigma_{\ell+1}(j')=1$ and $\mc R_{\sigma_{\ell+1}}$ maps the ordered commutator to the same unordered commutator are
\begin{align*}
    q_{\ell+1,\lceil \tau \rceil}!, \ldots, q_{\ell+1,1}!.
\end{align*}
Therefore, the combinatorial coefficient of $[\mc K_{v_{\ell+1,1}}, \ldots, \mc K_{v_{\ell+1,j'-1}},\mc K_{\lceil\tau\rceil},\mc K_{v_{\ell+1,j'}}, \ldots, \mc K_{v_{\ell+1,q_{\ell+1}-1}}]$ in the final sum is
\begin{align*}
    \frac{(q_{\ell+1}-1)!}{q_{\ell+1,\lceil \tau \rceil}!, \ldots, q_{\ell+1,1}!}q_{\ell+1,\lceil \tau \rceil}!, \ldots, q_{\ell+1,1}! = (q_{\ell+1}-1)!.
\end{align*}
In conclusion, we have
\begin{align*}
    \|\mc X_\ell(\tau)\|&\le \begin{multlined}[t]
        \bigg(\prod_{i=1}^{\ell+1} \frac{1}{q_i^2}\frac{1}{q_i!}q_i!\bigg)q_{\ell+1}/q_{\ell+1}\sum_{j'=1}^{q_{\ell+1}} \sum_{\substack{
                 v_{1,1}, \ldots, v_{\ell, q_\ell} = 1 \\
                 v_{\ell+1, 1}, \ldots, v_{\ell+1, q_{\ell+1}-1} = 1
            }}^{\lceil\tau\rceil} \bigg\| \Big[ 
                [\mc K_{v_{1,1}}, \ldots, \mc K_{v_{1,q_1}}], \ldots, \\ 
                    [\mc K_{v_{\ell,1}}, \ldots, \mc K_{v_{\ell,q_\ell}}], 
                    [\mc K_{v_{\ell+1,1}}, \ldots, \mc K_{v_{\ell+1,j'-1}},\mc K_{\lceil\tau\rceil},\mc K_{v_{\ell+1,j'}}, \ldots, \mc K_{v_{\ell+1,q_{\ell+1}-1}}]
                \Big] \bigg\|.
    \end{multlined}\\
    &=\begin{multlined}[t]\bigg(\prod_{i=1}^{\ell+1} \frac{1}{q_i^2}\bigg)
        \sum_{j'=1}^{q_{\ell+1}} \sum_{\substack{
                 v_{1,1}, \ldots, v_{\ell, q_\ell} = 1 \\
                 v_{\ell+1, 1}, \ldots, v_{\ell+1, q_{\ell+1}-1} = 1
            }}^{\lceil\tau\rceil} \bigg\| \Big[ 
                [\mc K_{v_{1,1}}, \ldots, \mc K_{v_{1,q_1}}], \ldots, \\ 
                    [\mc K_{v_{\ell,1}}, \ldots, \mc K_{v_{\ell,q_\ell}}], 
                    [\mc K_{v_{\ell+1,1}}, \ldots, \mc K_{v_{\ell+1,j'-1}},\mc K_{\lceil\tau\rceil},\mc K_{v_{\ell+1,j'}}, \ldots, \mc K_{v_{\ell+1,q_{\ell+1}-1}}]
                \Big] \bigg\|.
    \end{multlined}
\end{align*}
Finally, we integrate both sides from $\tau=0$ to $M$.  Since the right-hand side is piecewise constant in $\tau$, the integral is discretized to sum over $[M]$, where $\lceil\tau\rceil$ corresponds to the index $v_{\ell+1, j'}$. This yields 
\begin{align*}
    \int_{0}^{M} \|\mc{X}_{\ell}(\tau)\| \,\d\tau 
     \le & ~ \bigg(\prod_{i=1}^{\ell+1} \frac{1}{q_i^2}\bigg)\begin{multlined}[t]
    \sum_{j'=1}^{q_{\ell+1}}\sum_{v_{\ell+1,j'}=1}^{M} 
    \sum_{\substack{
                 v_{1,1}, \ldots, v_{\ell, q_\ell} = 1 \\
                 v_{\ell+1, k} = 1 ~(\forall k \neq j')
            }}^{v_{\ell+1,j'}}
     \bigg\| \Big[ 
         [\mc K_{v_{1,1}}, \ldots, \mc K_{v_{1,q_1}}], \ldots, \\ 
        [\mc K_{v_{\ell+1,1}}, \ldots, \mc K_{v_{\ell+1,q_{\ell+1}}}]
     \Big] \bigg\|, 
    \end{multlined} \\
    \le & ~ \bigg(\prod_{i=1}^{\ell+1} \frac{1}{q_i^2}\bigg) q_{\ell+1}
 \sum_{v_{1,1}, \ldots, v_{\ell+1, q_{\ell+1}} = 1}^M \bigg\| \Big[ 
         [\mc K_{v_{1,1}}, \ldots, \mc K_{v_{1,q_1}}], \ldots, 
        [\mc K_{v_{\ell+1,1}}, \ldots, \mc K_{v_{\ell+1,q_{\ell+1}}}]
     \Big] \bigg\|.
\end{align*} The right-hand side is precisely the definition of $\bar{\beta}_{\comm}^{(q_1, \ldots, q_{\ell+1})}$ from Eq.~\eq{def-beta-bar-general}, which concludes the proof.
\end{proof}
\end{document}
